% =====================================================================
%  GENERATED by ebook-method/etiske-undersoegelser/tr_parts.py build etiske-undersoegelser --write; do not edit by hand.
%  English translation of Harald Høffding, Etiske Undersøgelser (Kjøbenhavn: P. G. Philipsens Forlag, 1891), made
%  2026-10-16 from texts/hoeffding/etiske-undersoegelser/transcription.tex by four translators (TRANSLATOR-BRIEF.md,
%  GLOSSARY.md); Høffding's own English of two treatises (The Monist 1891, IJE 1890) used for the terms only.
%  \opage{N} = printed page N of the Danish begins.  Footnote marks *) ... †††) as in the print.
% =====================================================================
\documentclass[12pt,a4paper,openany]{book}
\usepackage[utf8]{inputenc}
\usepackage[T1]{fontenc}
\usepackage{libertinus}
\usepackage{libertinust1math}
\usepackage[german,french,english]{babel}
\usepackage{textalpha}
\usepackage[protrusion=true,expansion=true]{microtype}
\usepackage{setspace}
\usepackage[margin=1.2in]{geometry}
\usepackage{fancyhdr}
\usepackage{marginnote}
\usepackage{graphicx}
\usepackage{hyperref}
\hypersetup{pdftitle={Ethical Investigations}, pdfauthor={Harald Høffding}, hidelinks}
\pagestyle{fancy}
\fancyhf{}
\fancyhead[LE,RO]{\thepage}
\fancyhead[C]{\textit{H. Høffding: Ethical Investigations (1891)}}
\fancypagestyle{plain}{\fancyhf{}\fancyhead[LE,RO]{\thepage}}
\setstretch{1.3}
\DeclareUnicodeCharacter{0254}{\reflectbox{c}}   % ɔ (id est)
\renewcommand{\thefootnote}{%
  \ifcase\value{footnote}\or *)\or **)\or ***)\or \dag)\or \dag\dag)\or \dag\dag\dag)\else \arabic{footnote})\fi}
\makeatletter\@addtoreset{footnote}{page}\makeatother
\newcommand{\opage}[1]{\marginnote{\footnotesize\textit{[#1]}}}
\newcommand{\essay}[2][]{\chapter*{#2}\addcontentsline{toc}{chapter}{#1}}
\pdfstringdefDisableCommands{\renewcommand{\footnote}[2][]{}\renewcommand{\opage}[1]{}}

\begin{document}
\frontmatter
\begin{titlepage}\centering\vspace*{2cm}
{\large\textsc{Harald Høffding}}\\[1ex]\rule{3em}{0.4pt}\\[3ex]{\Huge ETHICAL INVESTIGATIONS}\\[3ex]
\fbox{\begin{tabular}{c}\small ON THE POSSIBILITY OF A PHILOSOPHICAL ETHICS.\\\small ON THE PRINCIPLE OF WELFARE.\\\small THE LAW OF RELATIVITY IN ETHICS.\end{tabular}}\\[2ex]
{\small [Etiske Undersøgelser]}\par\vfill
{\small COPENHAGEN}\\[0.5ex]{\small P. G. PHILIPSENS FORLAG}\\[1ex]1891\\[3em]
{\small Translated from the Danish}
\end{titlepage}
\chapter*{PREFACE.}
Of the three treatises that make up this work, I myself lay the greatest weight on the third, in which I have sought
to carry through an idea suggested in my Ethics (somewhat more in the German than in the Danish edition), which might be of some significance
for ethical thought. This treatise has already been printed in the first number of the ``International Journal of Ethics''
(Autumn 1890). --- In the two first treatises I have taken up for renewed treatment some questions that concern the first principles of ethics.
They may be regarded either as an introduction to or as a supplement to my Ethics (especially ch.~I, III and
VII). The second treatise (On the Principle of Welfare) is being printed at about the same time in the American journal ``The
Monist'' as part of a longer article. ---

\noindent\hspace*{2em}\textit{January 1891.}\hfill\textbf{Harald Høffding.}
\chapter*{CONTENTS.}
\begin{flushright}\small\textbf{Page}\end{flushright}
\noindent\begin{tabular}{@{}p{0.86\textwidth}@{\hspace{0.5em}}r@{}}
On the Possibility of a Philosophical Ethics\dotfill & 1\\
On the Principle of Welfare\dotfill & \textbf{24}\\
The Law of Relativity in Ethics\dotfill & \textbf{42}\\
\end{tabular}
\mainmatter
% --- p. 1 ---
\setcounter{footnote}{0}%
\essay[On the Possibility of a Philosophical Ethics]{On the Possibility of a Philosophical Ethics.}
\opage{1}
1. To preach morality is easy, to ground morality is
difficult, a philosopher has said. This proposition is correct,
if by preaching morality one understands merely the voicing
of the ideas of good and evil current in the given society.
But if by preaching morality one understands addressing the
individuals one by one in such a way that their character
thereby receives motives and incentives to develop not merely
in the spirit of the current ethical ideas, but in the spirit of
true and valid ideals, then it is not only the most difficult
and the highest art of all, but it will also necessarily lead
over to what that philosopher recognized as a great difficulty,
namely to give a grounding of ethical judgments.

Ethical judgments, judgments of good and evil, arise at
first as involuntary outbursts of feeling. They express
the value that is attributed to human actions or
institutions. In every society there is a certain sum total
of ethical judgments whose validity is recognized. They lie
in the air, in that the individual appropriates and applies
these judgments as immediately given and justified, just
as he inherits the instincts and traditions of the race in
general. This sum total of ethical judgments expresses what
has been called \emph{positive morality}. This has,
% --- p. 2 ---
\setcounter{footnote}{0}%
\opage{2}like positive religion and positive legislation, a historical
character. It has developed and continues to develop
constantly in reciprocal action with the fate and circumstances
of life of the race or society. Its content is a circle of
models, rules of life and commands, which for long ages pass
from generation to generation as ``unwritten laws'' and often stand
before the individual with a mysterious authority. The
individual submits to them all the more as they are
proclaimed and maintained by authorities, personalities,
whose physical or mental power awakens his fear, his
reverence, his trust or his love, and in the face of
which no urge to criticism is felt, or at any rate no
criticism easily gets a hearing. Such authorities are
parents, teachers, priests, princes, lawgivers, founders of religions.
--- Positive morality might be conceived as presented
in its connection. We should then get a presentation
of the ethical judgments valid among a certain people and at
a certain time. Theological ethics and the positive science of law
would be parts of such a positive\footnote[1]{The word ``positive'' denotes, in this whole context (positive morality, positive religion, positive legislation, positive ethics), not the opposite of ``negative''. Nor is it used in the sense in which Auguste Comte takes it, when he calls his philosophy ``la philosophie positive'' or Positivism. In Comte ``positif'' is the same as factual, really given in nature. Here, on the other hand, the word ``positive'' denotes what is set in a definite \emph{historical} form, as an expressly established determination or \emph{tradition}.} ethics (so
one might perhaps call the scientific presentation
of positive morality), in that theological ethics
concerns only the ethical judgments that can be derived from a
certain religious faith, and the positive science of law embraces only
the rules and commands that are expressly formulated
in the juridical laws of the land. Positive ethics would also
% --- p. 3 ---
\setcounter{footnote}{0}%
\opage{3}contain the rules and commands which public opinion,
custom, usage, tradition and instruction inculcate.

But when the fear of, or the trust in, the historically
given authorities wanes, or when for other causes
the urge to a more penetrating ethical knowledge awakens,
one cannot remain with positive morality.
A question arises about the very justification of the authorities,
about the first principles that condition the validity
of all ethical judgments. Do ethical judgments have no
validity, except insofar as they are maintained by authorities?
That is the great question by whose posing \emph{philosophical
ethics} comes into being. Its task must be
to show that, without presupposing the absolute
validity of the authorities, one can find a basis for ethical judgments
in human nature. Just as the philosophical theory of knowledge
concerns the first principles of the theoretical
conception of the world, which is grounded in human
nature, so philosophical ethics concerns
the first principles of the \emph{valuation} of human
actions, which is grounded in human nature itself.

2. The relation between philosophical ethics and
positive morality need not be hostile. Human
nature, in which philosophy seeks a basis
for ethical judgments, also reveals itself in history
and gains a determining influence on the development of
positive morality. The ideals and commands that make up the latter's
content have not arisen except by virtue of human
nature, as it acted under the given conditions. Even in the strictest
faith in authority, human nature therefore shows itself. No authority
can be purely external: authority is authority only through
the feeling it is able to awaken. Faith in authority
therefore always hangs together with the whole
% --- p. 4 ---
\setcounter{footnote}{0}%
\opage{4}life of feeling of men and through this again with the experiences that
men have in their struggle for existence.
By thus going back to the connection of positive morality
with human nature, one sees its kinship with philosophical ethics.
In part philosophical ethics has only the significance \emph{of bringing
the connection of our ethical judgments with our nature and
our conditions of life to full consciousness}, as indeed
philosophy generally can be said to be the self-knowledge of
the human mind in theoretical and practical respect.

But besides this relation between philosophical ethics and
positive morality, which comes about because the latter
does not have the same urge as the former to go back to the
first basis, three different relations may still obtain
between them.

Starting from the clearer idea
of the first principles, philosophical ethics may be led to demand that
work be \emph{carried further} essentially in the same spirit as
within positive morality, but from more comprehensive
points of view. Thus the feeling of humanity of more recent times
carries on the work on the basis of the principle of love of mankind,
which Christian ethics in particular maintained,
which Greek ethics made valid within
the limits set by national feeling, and which was already
at work in the ancient Indian ethics of the family, in the feeling
that embraced all who gathered around the same hearth. What
is at issue here is merely the \emph{extension} of the principle beyond the limits
of the confession, the nation and the family\footnote[1]{Cf. my ``Etik''. Copenhagen 1887. p. 326---331.}.

While what is at issue here is extension and carrying through,
it may in other cases be necessary to introduce an
entirely different motivation and grounding, in order that recognition
% --- p. 5 ---
\setcounter{footnote}{0}%
\opage{5}and continued application of the historically given may
become possible. The historically given in manner of acting or in the
ordering of social relations is then recognized and made use of
not from the motive that originally called it forth, but
from an entirely different motive. In other words, there takes place a
displacement of motives. --- From the sanctity of the hearth there developed
among the ancient Indians the right of asylum, which
became of special importance at the hearths in the temple and in
the king's court\footnote[1]{Cf. \emph{Leist}: Alt-arischen Jus gentium. Leipzig 1889. p. 20 f. 87. --- F. \emph{Pollock}: The King’s Peace (Oxford Lectures. London 1890). p. 70. 77.}; this right of asylum became in turn a form
under which peace could be established and pardon exercised;
and from this there sprang partly the state power's maintenance of
law and order, partly also the more humane treatment of
criminals and in general of those who are in conflict with
existing society. --- Another example is private
property, which has developed historically under the influence of
gainful activity, delight in activity and feeling of power, but
which society recognized and regulated in detail
for the sake of peace and order, and which is now recognized in ethical systems and
systems of political economy only insofar as
it promotes the individual's self-activity in the service of society\footnote[2]{Cf. my ``Etik''. p. 264---266.}. --- The displacement of motives has an illuminating
parallel in the way in which one and the same
organ in different species of living beings can perform
different functions, because it is adapted to different
conditions. The forelimbs are thus instruments of grasping
and of feeling in man, instruments of walking in the horse, instruments of flight in the bird, an instrument of swimming
in the whale. One and the same basic form appears transformed
in different ways under the different circumstances of life.
% --- p. 6 ---
\setcounter{footnote}{0}%
\opage{6}Conversely, different organs may come to perform the
same function, when living beings of different species
have in part the same conditions of life. Thus it is
highly different parts of the organism that function as organs of
respiration in mammals, fishes, crustaceans and insects;
the wings of birds and insects are due to the development
of entirely different parts of the body\footnote[1]{Cf. \emph{Gervais}: Élements de zoologie. Paris 1871. p. 49 f. 54 ff. \emph{Darwin}: Om Arternes Oprindelse. p. 507. 519. \emph{Körner}: Die logischen Grundlagen der Klassifikation der Organismen. (Wundt’s Philosophische Studien. II.). p. 242.}. The displacement of motives
in the psychological-ethical domain corresponds to
this transformation and adaptation in the biological domain.
In the biological domain too it is, in a certain
sense, the motive that changes, in that the same organs
come to serve different ends or different
organs to serve the same end. In ethics, actions
and institutions can thus acquire a different significance
and therefore a differently motivated recognition
than before. Just as little as one dares to infer from the ``adaptive''
(arisen by adaptation) similarity between butterflies
and true birds a real kinship
between them, just as little can one infer from the present ethical
significance of an action or an institution
to its historical origin and development. And just so
conversely: from the historical origin one cannot without more
infer to the value. --- The phenomenon of the displacement of motives
is therefore well suited to illuminate the relation between
the historical and the philosophical treatment of ethics\footnote[2]{See on this my ``Etik'' p. 8---10.}.

Finally, however, the relation may also be that the
principles and commands of positive morality must be rejected entirely,
and new principles set up in their place. The
content is then negative, not grounding,
% --- p. 7 ---
\setcounter{footnote}{0}%
\opage{7}continuing or transforming. Such an attitude is taken by
philosophical ethics to asceticism (insofar as this declares
feeling of pleasure in and for itself evil and wants to be more
than practice in necessary self-control), to intolerance,
to slavery and to other barbaric institutions. ---

3. We now return to the question whether there is
in human nature a firm basis for
ethical valuation, independent of the influence of the historically given authorities.
On the answer to this question depends
the possibility of a philosophical ethics.

All science has in the end its basis in human
nature. Thus logic presupposes the faculty of thought
(the faculty of comparison, the faculty of finding likeness
and difference); geometry presupposes in addition the faculty of
spatial intuition; physics and chemistry presuppose in addition
the various external senses (touch, sight,
hearing, taste, smell), etc. What now does ethics presuppose
for its part?

In ethical judgments there is an element at work that
plays no role in purely logical or natural-scientific
judgments. What is at issue in them is more than
establishing identity and difference or demonstrating a connection
between phenomena. The nature of an action
or an institution can call forth an ethical judgment,
a judgment of good and evil, only when the action
or institution in question awakens in us a feeling of pleasure or displeasure.
If we were merely intuiting, sensing and
comparing beings, without the capacity to feel pleasure and
displeasure, such judgments would not arise. Every judgment
arises through analysis of a given total conception\footnote[1]{See my ``Formel Logik''. 2nd ed. p. 13.}. The
peculiarity of ethical judgments is that they analyze,
and express in conscious form, the feeling that the sight of, or the
% --- p. 8 ---
\setcounter{footnote}{0}%
\opage{8}idea of, an action awakens in us.
The capacity to apprehend the action and to make its
character distinct to oneself (by comparison) is presupposed; but to this there is added
the element of feeling that it awakens, as a third factor
alongside knowledge of the actual circumstances and
the capacity to think about the action. Experience indeed also
proves that ethical discussions can be brought to a halt
in three ways, namely not only when one of the parties is convinced of a self-contradiction or of
a factual error,\footnote[1]{See ``Formel Logik''. p. 24 f.} but also when it is entirely different
kinds of feeling that determine the ethical judgments of the contending parties.
The ethical disagreements that concern only consequences
or factual circumstances are not of a decisive nature.
It is otherwise where it is a different feeling that determines
the valuation, --- where the \emph{motive of valuation} is thus
different.

If now feelings varied without limit, ---
if there were not certain typical forms of feeling,
which assert themselves with a certain firmness in the individuals
each by himself and in the same way in several
different individuals, then no philosophical ethics would be
possible. But it is not so. We find in every
normal individual a feeling that is the dominant one in
him, at least during a certain period, so that it constitutes
the center of gravity of his nature and will determine his
evaluative judgments. It is the most important constituent of
his ``real self''\footnote[2]{See my ``Psykologi''. ch. V, 5; VI E, 2; VII B, 1 c.}. And in every individual there will
also --- by means of memory and comparison, indeed
already by the mere instinct of self-preservation ---
assert itself a connection between the single
% --- p. 9 ---
\setcounter{footnote}{0}%
\opage{9}moments and the single periods of life. In this the ``formal self'' manifests itself\footnote[1]{``Psykologi''. ch. V, 5 (cf. IV, 7 e).}.
There will therefore manifest itself a stronger
or weaker tendency to evaluate actions not
merely by their influence on the single moment or
the single period, but by their influence on life
as a whole. And in the different individuals there manifest themselves
not merely feelings that are attached to their existence
each for itself alone, but also such as are determined
by the common conditions of life, and which immediately link the
one individual to one or more others. The maternal instinct and
the feelings springing from it are immediate natural promptings
in this direction. Within the family, the tribe
and the nation, and within the various associations
and fellowships, more or less kindred
feelings develop under the influence of
common work and common fate. The intensity and
extent of such feelings stand as a rule in inverse proportion to each other;
they stir most strongly and powerfully in the narrowest circles,
especially in the family and the neighborhood. But history shows that
they have been able to develop so as to embrace the whole of mankind.
The idea of a mankind and of universal
love of mankind (caritas generis humani) developed
already among the Greeks and Romans of antiquity in the last
centuries before the birth of Christ, probably as a consequence of
the way in which the national barriers were broken down
by the conquests of Alexander the Great and of the Romans, so
that the feeling of kinship and community became possible.

The ethical judgments must now be determined by one of the
feelings here mentioned (besides being determined, as already stated,
by the consistency of thought and knowledge of
the actual circumstances). But therein lies that \emph{there is a
possibility of as many ethical systems}
% --- p. 10 ---
\setcounter{footnote}{0}%
\opage{10}\emph{as there are different dominant feelings that can
determine the judgments}. The different ethical
systems differ from one another by the motive of valuation.
The history of ethics shows us that the different possibilities
have to a great extent been made use of. Even
a pure ethics of the moment, which lets regard for the single
moment, abstracted from all other moments, be sole ruler,
has been tried (by Aristippus of Cyrene). One has
often taken as the point of departure the self-preservation, enjoyment of life
and feeling of power of the single, isolated individual. Thereby
one is led to a system in which regard for other individuals
acquires value only as an often indispensable means
for attaining goods for the individual himself. Ethics thus becomes a
doctrine of the necessary detours the individual must take
to attain his happiness. Where basis and motive of valuation
are taken from sympathy, different systems may again
become possible, according to the extent sympathy has.
An ethics of the family is approximately realized in the
ancient Indian ethics. A national ethics was developed by
Plato and Aristotle. An ethics of estates or of associations
appears approximately in history, also in our days,
in the mutual conflicts of the different classes of society.
But an ethics of humanity too works its way forward
through history; the valuation of actions and institutions
then takes their influence on all human
beings into consideration.

\emph{Psychologically and historically viewed} there is a continuous
connection between all these different standpoints,
which in theory and practice appear carried through
with more or less consistency and with greater or
lesser expertise. The connection rests on the fact that the
one motive of valuation can, by virtue of psychological laws and
under the influence of the experiences of life, be replaced by the
other. With such a change the
% --- p. 11 ---
\setcounter{footnote}{0}%
\opage{11}system of ethical judgments also changes; but this change is
effect, not cause. It is the arising of new feelings
that matters, and they arise through practical experiences,
not through theoretical reasonings. A practical
education now takes place constantly, not only within the family,
but also within social life as a whole. The
individual learns to subordinate regard for the single
moment to regard for the whole of life, and regard
for his self-preservation and enjoyment to regard for
what the common, general good requires. An attachment to
other individuals, initiated by egoistic motives,
can under continued life together and its experiences become an
independent object of feeling of pleasure; what was at first
grasped as means can become end. In this kind of displacement of motives
consists to a great extent the education of mankind
through history. But it would hardly go far,
if it were not supported by immediate sympathetic
instincts\footnote[1]{Cf. ``Psykologi'', VI C, 2---7.}.

\emph{Ethically viewed}, on the other hand, there is discontinuity between the
different standpoints and systems, since it is not
possible for the one to refute the other by way of logical
proof or by means of objective facts.
The greater the consistency the single system possesses,
the more unassailable it will be. The resolute and
consistent egoist, who makes himself the sole end
and regards other individuals merely as means for himself,
cannot be refuted from the standpoint of the ethics of the family, the ethics of the nation
or the ethics of humanity. He says with Shakespeare's
Richard III: I am myself alone! Naturally it is
only insofar as the system is consistent that it is
unassailable; if family feeling, national feeling or
love of mankind stirs in the egoist, he will
% --- p. 12 ---
\setcounter{footnote}{0}%
\opage{12}be able to come into contradiction with himself and be easily refuted.
Otherwise he can be brought to give up his standpoint only
by a feeling being awakened in him that he has hitherto not
known. But this happens by ``moral preaching'' (in the highest
sense), not by mere ``moral grounding''. It is ethical
art, not ethical science, that then goes
into action.

There are in every society individuals who do not stand
on the basis that is necessary in order that social life
at the given stage may subsist. Such individuals
are children, the insane and criminals. One may say that
they have other ethical systems than the one society as
such holds. Toward them other means are applied
than mere reasoning. They become the object of
a practical mode of treatment, a discipline and education that
is to awaken in them the feelings they have hitherto lacked.
Society arrogates to itself, \emph{by virtue of its own ethics}, the right
to treat them in this way, but cannot
always ground this right on the ethics of such individuals themselves.
--- In every ethical judgment we utter there is really
such a pedagogical moment. The ethical judgment springs
not merely from feeling; it also speaks to feeling.
Him whose action we evaluate, we praise or blame;
we thereby set his feeling in motion, make ourselves an
authority over him and bear responsibility for the effect
of the motives we can in this way call forth.

4. Meanwhile it must be well noted that the different
systems are not purely coordinate with one another. The single
system can give the other systems a certain subordinate
significance within itself. The system of the egoist, e.g., does
not exclude regard for family and state, but regards
them merely from the point of view of means. The ethics of the nation and of humanity
must from their standpoints show how
the single individual within their systems comes to
% --- p. 13 ---
\setcounter{footnote}{0}%
\opage{13}his full right, though in another way than according to the egoistic
system. Through sympathy, through devotion
to a larger whole and work in its service the individual's
powers and abilities are developed; his personality is widened
and strengthened precisely through self-forgetfulness. To
forget oneself is not the same as to lose oneself.
It is to set oneself other ends than one's own ego.
Through the increased and heightened striving that these more far-reaching
ends condition, the mind gets greater life, strength
and flight. It is thus an enrichment, a gain for
the individual, if sympathy gains power in him. But this
can be experienced only when sympathy \emph{has superseded} egoism,
and thus cannot be ground or motive for this
supersession. Therefore sympathy would not have developed
in the human race, if nature had not here, just
as in self-preservation, laid hold on us through instincts
and involuntary experiences. We should otherwise
stand to all eternity like the man who would not go into the water until he
could swim.

One might perhaps think that things went here as in
a physical or chemical experiment, by which something is established
by way of experience (a posteriori) that could not
be established by way of mere thinking (a priori),
and that ethics was thus here not placed otherwise
than any science that takes its presuppositions from experience.
But one would then overlook a very essential
difference. The experiment in question is to take place
with the individual himself, and precisely with that in him that
conditions his ethical judgment. When the individual has experienced
how sympathy widens his ego, makes it healthier,
stronger and opens to it new sources of satisfaction,
then he has himself become another than he was before; he
has acquired another real self. He draws up the account
from another basis than that on which he stood before he
% --- p. 14 ---
\setcounter{footnote}{0}%
\opage{14}undertook the experiment, --- and the reckoning is determined throughout by
the new basis, the new motive of valuation.
He thus does not change his view by inference from
the experiment, but the latter immediately changes his feeling
and as a consequence his view.

Moreover the transition from the one basis to
the other does not take place with the same ease as that with which
an experiment is performed. History and daily
experience show this sufficiently. The education of mankind takes
a long time. For some individuals the possibility has
even been denied that the minimum of sympathy that present civil society
requires can be awakened; in another connection we shall discuss
this denial more closely. Experience shows in any case,
that there are individuals who, in consequence of the
basis on which they now once stand, collide with the given
conditions and perish. Tragedy
draws some of its greatest subjects from the clash
of ethical systems; within one and the same system tragic
clash is not so easily possible. --- And even where the transition
from one basis to another succeeds, it as a rule takes time and takes place
often through severe crises.

It would therefore be a very inexpedient use
of words, if one were to call a feeling egoistic,
because it is after all in sympathy too our \emph{own} pleasure or
displeasure that we have, and because sympathy could not subsist,
if it did not bring with it a fuller and more inward life
and thereby a greater satisfaction and self-assertion than
would be attained, if feeling of pleasure and displeasure were determined only
by a self-preservation and self-assertion for which everything
else in the world were mere means. Such a use of
words would mislead one into overlooking some of the most significant
and most deeply penetrating oppositions in the domain of the
life of the mind.

% --- p. 15 ---
\setcounter{footnote}{0}%
\opage{15}5. It rests on a choice what standpoint is taken
in a presentation of ethics. It would be of great interest to have the
various possible systems each by itself fully worked out,
so that we could see how ethics shapes itself in detail
at a standpoint of the moment, an egoistic standpoint, a family standpoint, etc. It
would throw light on the nature of the motives of valuation. The interest of
historical or comparative ethics rests on this; for historically we
find approximations to the various possible systems.
But also only approximations. For in real life the motives of valuation
(just as well as the motives of action) are always more or less
composite. Our conscience (that is to say, the feeling
that actually determines our valuation) hardly ever coincides
completely with sympathy or egoism or with
a single one of their forms. Feelings of highly different kinds
work together in it. When ethics, in order to reach full clarity
and consistency, picks out a single definite feeling
and holds to it as basis, in that it conceives
it in full purity and strength and with greatest extent,
it performs a similar abstraction and idealization
as every formal science performs in its first
step. The pure space which geometry presupposes in its
first definitions and axioms does not occur in reality
as an object of sensation and intuition.

Now when ethics chooses as basis an ideal conscience,
which in its valuation is determined by universal
sympathy, it has indeed performed an abstraction, but this
consists in its conceiving in ideal perfection a force
that actually exists in the world, has worked its way forward
and probably will continue to work its way forward. It has
thus in spite of its ideality firm ground under its feet.
It proposes to us to regard the human world
and what goes on there from this vantage-point and see what
% --- p. 16 ---
\setcounter{footnote}{0}%
\opage{16}ideals and tasks follow from it. Others may then regard the world from their vantage-point. When we then each for himself
proceed with as great consistency as possible and rely
on as many and as certain experiences as possible, our
knowledge will have made the progress it is able to make in this domain.
The possibility of grounding ethics
leads us to the necessity of recognizing a limitation
of its scientific character. Grounding and limitation
belong closely together in all domains. ---

When sympathy is chosen as motive of valuation,
it lies therein that only the ends which an enlightened, consistent
and universal sympathy acknowledges can be
the highest ethical ends. But it does not lie therein that
sympathy necessarily ought to be motive of action in
all cases and in all individuals. There may lie in human
nature other motives, which need not
lead to conflict with the ends sympathy sets, but
may precisely work for their benefit; and these motives may even
be indispensable. Urge to self-preservation
and self-assertion, ambition, adherence to custom and
usage or to venerable authorities --- these are examples
of forces and tendencies in human
nature which an ethics that rests on sympathy as
ultimate motive of valuation can fully acknowledge, insofar as
they act in the direction of the ends it has set.
The difference between such motives and universal
sympathy is that, whereas the latter \emph{necessarily}
comes into contradiction with itself by not regarding
actions and institutions from the standpoint of the whole (of society,
ultimately of mankind), the other motives do not necessarily come into
contradiction with themselves by a more limited range of vision. For the egoist
it is indifferent how it goes with the others (whether
they too are egoists or not), when they only promote
% --- p. 17 ---
\setcounter{footnote}{0}%
\opage{17}\emph{his} ends. And even if, as some have thought they could
prove, there should be a harmony among all
egoists' well-understood interests, then the single
egoist rejoices over this only insofar as this harmony of ends benefits
\emph{him}. If the idea of such a real or possible harmony
of interests aroused in him a joy or perhaps even enthusiasm,
which goes beyond the joy that is felt over a means
that has worth only as means, then he would no longer
be an egoist; he would then have reached the point of having other
ends besides himself.

The ethics grounded on sympathy as motive of valuation
leads each individual to regard himself not merely as
end, but also as means, and to regard others
not merely as means, but also as ends. From no
other basis can this way of regarding be grounded.
Its general basic thought can be expressed thus:

Look at your life and at the lives of the others not merely from
your own vantage-point, as if it were the center of the world, but
from a higher vantage-point: that of society, of the race! Because you
reach the goal you have set yourself, it is not said that all
others too have reached their goal. Because you have had
your wishes fulfilled, there is suffering enough still left in
the world, justified but unfulfilled wishes in sufficient
quantity. Were you even the greatest that has lived, you are
yet only one among many of the race whose history
is not concluded with your history!

6. An essential objection to the attempt to
build ethics on sympathy as motive of valuation might
seem to lie in the struggle for existence as the guiding
point of view for all life, human life too. When
life is an uninterrupted struggle to subsist between different
species and between different individuals of the same
species, does it not then hold for each individual to get ahead of the
% --- p. 18 ---
\setcounter{footnote}{0}%
\opage{18}others and to procure for himself as great advantage as
possible at their expense, and must not an ethics grounded on sympathy
then lack all root in real life?
There seems indeed then to hold no other opposition than
that which exists between those who can subsist and
develop, and those who cannot, --- between the
sound, strong and happy on the one side, the sick,
weak and suffering on the other side. And is there here
not an absolute, a hopeless opposition, which
cannot be shaken without general enfeeblement and paralysis
becoming the consequence?

While some interpret the Darwinian doctrine of
the struggle for existence in this way and with aesthetic
complacency revel in the thought of the ruthlessness
and strength which the struggle for existence leads to unfolding, there are
others who indeed concede that such ethical consequences
lie in that doctrine, but who find therein one of the
most essential objections against it. --- As for Darwin
himself, he conceded none of these consequences.
As a moral philosopher he joined the ethical school founded by Shaftesbury
and Hutcheson and carried further in our days by Herbert
Spencer, which saw in sympathy the basis of
ethics, a conception which in this treatise it has been attempted to give a more definite illumination and
carrying through. I shall now try briefly to show
that the great naturalist had a more correct conception
of his doctrine than that which appears in those
who draw those supposed consequences.

7. The struggle for existence brings to unfolding all possibilities
that are of help under the given conditions.
There can thus develop a ruthless urge to self-preservation,
which can rise to different degrees of conscious
egoism. There lies a possibility and a tendency in the
human nature in this direction. But it is not
% --- p. 19 ---
\setcounter{footnote}{0}%
\opage{19}the only possibility that lies in human nature, nor
the only possibility that can unfold under the struggle
for existence. This struggle is not in the literal sense a
``war of all against all'', as Thomas Hobbes in his time
assumed. The existence at stake in the struggle is not merely
the single individual's, but the species', the race's. And the species is maintained
not always and not merely by single individuals
each by himself as ruthlessly as possible seeking to hold
themselves up, but often also, yes: only, by their striving and
fighting in common, or by the one individual sacrificing
himself for the other. The struggle for existence thus produces
not merely egoism, but also means of getting
beyond egoism. Its effects are not all unambiguous,
not all of the same value in ethical respect. An
ethics is precisely necessary because a choice must be made
between these effects, as soon as the question is
which of them we with our conscious willing and intervention
shall favor, and which we shall seek to check
or displace. As willing beings we are not merely
effects of the preceding development, but also
(in part at least) causes of the following development,
and our ethical deliberations concern the direction
in which our causality ought to go. To decide this,
we need a standard, an ideal, and the first task of ethics
is to establish the basis from which such a
standard can be found. Philosophical ethics seeks,
as we have seen, this basis in human
nature. The struggle for existence cannot, as certainly as it is
the condition of all life, be stopped; but it may
perhaps be conducted in a way that accords with the given
motive of valuation.

By looking at the first conditions of the preservation of the race
we shall find a possibility for this. --- Goethe
% --- p. 20 ---
\setcounter{footnote}{0}%
\opage{20}said, when he saw a sculpture representing a cow with a suckling calf: ``Here we have a subject of
the highest kind: the world-sustaining nourishing principle running through the whole
of nature meets us here in a beautiful simile''. The Darwinian doctrine is
a closer elaboration of this thought. It has shown that without instincts,
which lead to care and self-sacrifice for the offspring,
a species cannot subsist in the struggle for existence. However
many riddles still remain, as regards the first germs
of the most wonderful instinct in nature, the maternal instinct
or the parental instinct, its necessity is evident. It shows itself
especially where it stands in the most glaring opposition to a being's
other instincts. Spiders, e.g., are beasts of prey, which seek their prey each
by himself by means of their webs, without there, as in
so many other beasts of prey, being able to develop a certain social
feeling during the hunting-fellowship. They attack and devour
on the contrary other individuals of their own species, even
when there is no lack of other prey. Not even in
the mating season are male and female safe from each other's predatory appetite.
The offspring would therefore at once perish, if
not even in these otherwise so ``egoistic'' animals the maternal instinct
appeared even in unusually high degree\footnote[1]{See \emph{Fr. Dahl}: Versuch einer Darstellung der psychischen Vorgänge in den Spinnen. (Vierterljahrsschrift für wissenschaftliche Philosophie. IX). p. 180 ff.}. ---
The instincts and feelings that bind parents and children
to each other are thus decidedly favored by the struggle for
existence. Only in rarer cases are the opposite instincts
object of natural selection. Thus the instinct that makes the worker bees kill their
brothers, the drones, and the queen bee kill the
% --- p. 21 ---
\setcounter{footnote}{0}%
\opage{21}bee-princesses\footnote[1]{\emph{Darwin}: Menneskets Afstamning. I. p. 75. --- Arternes Oprindelse. p. 259.}, can be explained only by its significance in the struggle for existence,
which thus under some circumstances
favors maternal love and love of siblings, under
other circumstances the opposite; yet this last
is the rarest case. Where within one and the same species
single individuals in greater or lesser degree lack
the parental instinct,\footnote[2]{See on this \emph{Lombroso}: Der Verbrecher. Deutsche Uebers. p. 9. 27. --- Dictionnaire des sciences anthropologiques. Article ``Famille''.} such exceptions must be explained
in the same way, insofar as they are not remnants from earlier
stages of development, which maintain themselves in single cases
by atavism.

This way of regarding finds application also
to the significance of solidarity and community for the
single individuals in the maintenance and conduct of their lives.
Not as a consequence of conscious calculation and arbitrary
choice, but in consequence of the nature of the circumstances, comes
the single individual to carry on the struggle for existence
together with others. He practices solidarity
and cooperation before he can form any idea
of their significance. There forms hereby within one and
the same group of human beings a feeling of brotherhood and
humanity, which perhaps at first only
shows itself in war and danger, but later can express itself
also under other circumstances.

Under the common life thus involuntarily instituted and practiced,
under community in fate and work,
in pleasure and pain, the individual learns practically to
apply a higher standard than that which can be fetched from
his own purely individual interests. Through life together
and joint activity the fellow-feeling develops, ---
% --- p. 22 ---
\setcounter{footnote}{0}%
\opage{22}irrespective of what motives have originally led to the life together
and the joint activity, for a displacement of motives
often takes place here.

In my ``Etik'' I have shown how this practical
development of the feeling of community plays a role within
marriage (p. 180), in the organization of the working class
(p. 244), in the development of religious freedom (p. 312)
and in national feeling (p. 338 f.). I have traced it back
to ``the Aristotelian principle'' (p. 162 f.), that involuntary
practice of qualities and skills precedes
conscious possession and exercise. Qualities
form, --- that was Aristotle's doctrine, confirmed by more recent biology
and psychology, --- by working and acting, and where a quality or a skill has
formed, a feeling of pleasure is connected with the exercise
of the corresponding actions. The first condition
is that the circumstances of life give occasion for practice.

8. One more main point is to be emphasized to
illuminate the relation of ethics to the struggle for existence.
This concerns not merely the pure, bare existence, the
mere physical subsistence. The existence that is fought for
is that which is led at a definite stage of life. Not without
resistance, not without several possibilities first being tried,
does life go back to a lower, more elementary stage.
Every stage of life has its needs, in which its
essence reveals itself. Only at the very simplest
stages of life is the feeling of pleasure and displeasure attached
solely to what immediately serves to maintain life. When these
most elementary needs are satisfied, there will
at higher stages stir an urge to use the powers and a pleasure in this use, even if it does not
lead to any practical result. The feeling of pleasure in play
and the development of the aesthetic feeling are examples of this.
The joy in knowledge of truth is also an example of
% --- p. 23 ---
\setcounter{footnote}{0}%
\opage{23}this displacement of motives, by which the feeling of pleasure is attached
to the activity itself and not merely to its result.
A higher stage has been reached, where the elementary demands of life
recede in such a way that ideal interests, independent of
the physical demand of self-preservation, become possible. And
when such a stage has been reached, the struggle for
existence concerns the preservation of this stage of life: life is that which is
led at this stage; everything else is death. With right one therefore speaks
of a mental death, where ideal interests
are extinguished. To such ideal interests belong now
also sympathy and conscience, when they have worked
their way forward. A life that is won only by denying
them would now stand as worthless\footnote[1]{Cf. ``Om Grundlaget for den humane Etik''. (Copenhagen 1876), p. 30 f., where I have already put forward this thought.}.

% --- p. 24 ---
\setcounter{footnote}{0}%
\essay[On the Principle of Welfare]{On the Principle of Welfare.}
\opage{24}


1. When the principle of welfare is recognized as the \emph{standard}
in the valuation of actions and institutions,
this means that they are called good or evil according
as in their effects (as far as we can trace
them) they bring about a predominance of feelings of pleasure or a predominance
of feelings of displeasure in a larger or smaller circle of conscious
beings. Every action can be compared to a
stone thrown into the water; the movement propagates itself in
larger or smaller circles, and the valuation depends on
how it works there, where it intervenes in the weal and woe of living
beings. Just as theoretical science
explains the single natural phenomenon by its connection
with other natural phenomena, so ethics evaluates the
one feeling by its relation to other feelings: the
satisfaction the agent himself feels in his acting
can be approved only if the effects of the action do not
intervene disturbingly in the feeling of pleasure of other beings, when
such an intervention cannot be shown to be a necessary
means to a so much the greater feeling of pleasure, --- and the
best actions are those that increase the feeling of pleasure in the world
and diminish the feeling of displeasure. --- As a
\emph{standard}, this principle corresponds immediately to sympathy as \emph{motive of valuation}.
How far it is possible to recognize it
% --- p. 25 ---
\setcounter{footnote}{0}%
\opage{25}from standpoints that lay other motives of valuation
at the basis, shall not be discussed further here\footnote[1]{On this see ``Etik'' p. 39---42.}.

The valuation does not stop at the external action,
but goes back to the motives of the agent, to
his traits of character, to the whole inner life from
which the action has sprung. This has its ground
in the nature and origin of valuation. Ethical judgments are originally
involuntary outbursts of feeling. And these outbursts of feeling
are of great practical importance, in that they act
as determinants on the will (one's own and that of others) and produce
motives for future actions. It is therefore necessary
to go back to the point where it is possible to
influence the force that has produced the action, and
this point lies in the inner life, in the disposition of the
agent, in the motives that have brought about the action.
Only when they are changed are the future actions also changed,
and with these in turn their effects. It is
the tendency of effects that is consistently traced back to
its source.

A force we know only from its effects. Only
by inference can we form an idea of what
goes on in the mind of a human being, of the motives and
faculties he commands. A goodness or greatness that
never at any time showed itself in action could not
become the object of ethical recognition. It would rest
on self-deception. There must at least take place an
inner work, a yearning and striving in the direction the
ethical principle demands. Often the individual in his valuation
of himself must take refuge in this inner work,
and the true ethical valuation must in the end concern that\footnote[2]{On this see the following treatise on the law of relativity in ethics.};
but it is indeed precisely the beginning, the first attempts at that which
% --- p. 26 ---
\setcounter{footnote}{0}%
\opage{26}the principle of welfare demands, only that the continuation is often made impossible on
account of individual circumstances.

But just as important is it to maintain that the effect
of the action need not show itself at once.
When great and good men are so often misjudged by
their own contemporaries, this is explained by the fact that only a very
comprehensive view can measure the significance of their conduct.
Their greatness and goodness rest precisely on this, that their thought,
feeling and will embrace far more than their short-sighted
and narrowly feeling contemporaries are able to discover. A long
time may pass before general understanding and appropriation
become possible. Therefore it by no means lies in the principle of welfare
that one should in one's conduct adjust oneself wholly to
the want and the needs that men have at the moment.
The principle of welfare demands precisely that we not shrink from
the struggle against prejudice and inertia. The best one
can do for others is often simply to let them feel
that in their wishes and needs they still stand too
low and do not make sufficiently great demands. Thus
--- to take an example --- the great artist
often goes his lonely way, not understood or even misunderstood
by the great multitude. And yet he follows the principle of welfare,
without perhaps thinking of it, by rigorously
upholding the demands of art. He increases the mental
capital of the race and bestows on it a force which in the future
may come to act in large circles. Only a
short-sighted conception and application keeps to the
want prevailing at the moment instead of looking to
the permanent conditions of life of the race and to the steady sources
of new life and new activity\footnote[1]{This last argument is cited here from my ``Etik'' p. 94.}.

2. The principle of welfare furnishes a \emph{rule} we can
lay at the basis of our valuation of actions. But
% --- p. 27 ---
\setcounter{footnote}{0}%
\opage{27}it by no means demands, as has sometimes been supposed,
that regard for welfare shall also be the \emph{motive} for every
action. We resort to general principles only
in order to orient ourselves in such cases where the immediate
appraisal or the instinctive feeling cannot bring
the matter to a decision, that is, in cases of doubt, or else
when we strive for a systematic treatment of ethical
questions. The ethical feeling can work quite involuntarily
and without actual reasoning, in that we are immediately
captivated by the action as it stands before us (before
thought, that is as possible, or in experience, that is as
real), just as in our aesthetic feeling we can immediately
be gripped by the beauty of a work of art or of a landscape,
without needing to resort to aesthetic reasonings. Or
one follows with confidence the ``unwritten laws'' contained in
custom and usage, in tradition, in sum in positive
morality. Precisely by virtue of the principle of welfare such an
immediacy is to be recommended and maintained, as
long as it does not lead to overlooking real problems
and tasks. It is the state of innocence, from which no one
ought to be torn without necessity. When the immediate
security is broken, the abstract principles become
necessary aids; but they can often be applied to the single, concrete cases only through
many intermediate links,
and they do not so easily gain a practical influence
on the will.

Even quite other motives than ethical feeling or
devotion to the demands of positive morality can
the principle of welfare approve, wish or require. It is thus
most fortunate if a man takes care of his wife
and his children because he loves them, and not because his
ethical feeling requires it. Where duty has consciously to
come into force in the most intimate relations between human beings,
it is as a rule a sign that a less fortunate relation
% --- p. 28 ---
\setcounter{footnote}{0}%
\opage{28}is present. Perfect love drives out not only
fear, but also duty. Like the family relation,
the relation of the artist and the researcher to their
field is also an example of this\footnote[1]{Just as a distinction is made here between standard or principle of valuation and motive of action, so in the preceding treatise (§ 5) a distinction is made between motive of valuation and motive of action. Because of the close connection between principle of valuation (standard) and motive of valuation (basis) these two considerations nearly coincide.}.

Wundt\footnote[2]{Ethik. Eine Untersuchung der Thatsachen und Gesetze des sittlichen Lebens. Stuttgart 1886. p. 339.} has put forward the following objection
to the principle of welfare: ``It is conceivable that one person sacrifices
himself for another; it is conceivable that one surrenders property
and life for some ideal end, for the fatherland,
freedom, faith or science. But it has never
happened and will never happen that anyone renounces,
in order that the sum of happiness that is given in the world may become
greater.'' This objection overlooks precisely that what
is principle for the valuation of an action, when it is
recognized as good, does not \emph{need} to be motive for
the action. The thought and feeling of the agent may
very well stop at the fatherland, freedom, or whatever
other ideal end he sacrifices himself for, without his needing
to make any reflection on what conditions
the value of these ideal ends themselves. But in a systematic
ethics or in practical cases of doubt we discuss the worth and significance for human life
of love of country\footnote[3]{See my ``Etik''. ch. XXXVI.}, of the feeling of freedom\footnote[4]{Ibid. ch. XXIII and XL.}, of faith\footnote[5]{Ibid. ch. XXXI---XXXIII.} and of science\footnote[6]{Ibid. ch. XXVIII and XXIX.}.
If, e.g., freedom were not a good for a people, it would
% --- p. 29 ---
\setcounter{footnote}{0}%
\opage{29}be wrong for a person to be willing to sacrifice his life
to preserve it. Of course it is never in the single
case a question of the abstract concept of welfare in and
for itself, any more than in a special theoretical
investigation it is a question of the abstract concept of cause
in and for itself. But we set up in ethics the principle of welfare
and in the theory of knowledge the principle of causality, when we by way of
analysis have sought the first presuppositions for
our practical and theoretical (evaluating and explaining)
mental activity. ---

3. It does not conflict with the principle of welfare that pleasure
so often must be bought with pain. For the pain then stands
merely as the necessary transitional step, and the significance of the principle of welfare
is precisely to maintain that the burden of proof
rests upon him who asserts that such a transitional step
is necessary. Every infliction of pain must be motivated,
whereas the feeling of pleasure in and for itself (when that which calls
it forth does not at the same time cause greater unfortunate effects)
is justified. The principle of welfare says merely: produce as
much pleasure and as little displeasure as possible by your
acting! To what degree it is possible to fulfil
this demand, on that the principle of welfare in and
for itself contains nothing. A principle is not overthrown by the difficulty
of applying it. And experience shows that there is
a feeling of pleasure that is not bought too dearly, even if deep pain
goes along with the purchase. Against the principle of welfare
Clärchen's song (in Goethe's ``Egmont'') has been adduced:
``Himmelhoch jauchzend, \emph{zum Tode betrübt}''. But
let Clärchen speak out, and let us note the song's
last lines, in which she draws the result of her whole
series of moods: ``Glücklich allein ist die Seele, die
liebt!'' The point is: there is a movement in the mind, a
life in feeling, which is connected with so deep and great
a satisfaction that the strong oscillation between pleasure and
% --- p. 30 ---
\setcounter{footnote}{0}%
\opage{30}pain does not cancel the total feeling of happiness, but precisely
strengthens it. Two psychological laws work together here.
The one is that according to which the pain (or displeasure,
or sorrow), when it does not exceed a certain limit,
can form a background for the feeling of pleasure and thereby heighten
it. Already in this lies a sufficient motive for
choosing such states instead of those that do not reach
so high in intensity, although they are on the other hand
more uncompounded states of pleasure. The other law is that
in sorrow itself there can be something attractive, probably
because the strong life, the vigorous movement,
the tension of the faculties that goes with it, in and
for itself brings a satisfaction. All exertion of force
that is not disproportionate is connected with feeling of pleasure.
In the midst of sorrow a feeling of pleasure in the very
strength and violence of the sorrow can thus assert itself, which
shows itself in that one does not \emph{will} to be torn from it. The organic
discharge that is connected with every strong
mood is an important element here, even if it should not
contain the full explanation. --- When Auguste
Comte had lost the woman who exercised so decisive
an influence on the direction of his mind in the last
period of his life, he once said in an address
he directed to his memory-image of her: ``Only
to you do I owe that I shall not come to leave life
without having in a worthy manner come to know the best feeling of human
nature\dots{} In the midst of the
heaviest pains that the feeling can bring, I have never
ceased to feel that \emph{the most essential thing for happiness} is
to have the heart filled in a worthy manner, even if it
is with pain, indeed even with the bitterest pain.'' --- Auguste
Comte and Clärchen are thus quite in agreement, and the morality of welfare
is in agreement with them. With the great movement of life
and fullness of life there follow also the great costs, the great
% --- p. 31 ---
\setcounter{footnote}{0}%
\opage{31}contrasts and the great oscillations. It is a matter then of
having courage to venture and choose. Who has said that the
highest was to be had for a low price?\footnote[1]{This last argument I have already set forth in ``Tilskueren''. Nov.---Dec. 1889. p. 861 f.}

We simply have no other psychological criterion
of the health and vigor of life than the feeling of pleasure. That
which --- in \emph{all} its nearer or remoter effects ---
causes only feeling of pleasure in all the beings it affects,
cannot possibly be rejected. To this extent welfare,
defined as lasting feeling of pleasure, is the ultimate principle of valuation.
It does not conflict with this that the highest we
actually know often, through some of its conditions, entails
great pain. Whoever wants to be secure against privation and
sorrow must not love anything whatever; but if now
love itself were the highest happiness?

Pain is everywhere a sign of a beginning dissolution
of life\footnote[2]{On this theory see: ``Psykologi'' VI D, 2---3.}. In the simplest form this shows itself in
``physical'' pain, which arises by the tearing apart of the
organic tissue. It holds also for ``mental''
pain, which consists in sorrow, doubt or repentance\dots{} It
points to a disharmony between different forces and
impulses in the mind, a disharmony that, by rising, finally
can burst consciousness. When pain is nevertheless often
a necessary transitional step, it is partly because it
(as already shown) by contrast and stimulation can set
stronger forces in motion, which outweigh the dissolving
forces, partly also because it concerns the dissolution of something
in us that hinders a freer and richer unfolding of life. Childbirth
is connected with pain, because the new life can appear only
by detachment from the old, of which it has
hitherto formed a part. Knowledge of truth is won
often with pain, because prejudices and illusions burst.
% --- p. 32 ---
\setcounter{footnote}{0}%
\opage{32}In the pain of repentance a lower self is dissolved, so that a higher
self may develop.

4. What has especially occasioned the resistance to the principle of welfare,
is surely that one is inclined
at the word feeling of pleasure to think first and foremost of
the most elementary and sensuous forms of pleasure. The principle of welfare
does not exclude these, but includes more
than them. It takes account of all sides of the nature
of man, in that it maintains that a lasting feeling of pleasure cannot
be produced with certainty when an essential side of
this nature is left unheeded. The defect of the elementary
feelings of pleasure is that for the most part they correspond
only to momentary and limited conditions. A being whose
feeling is of a purely elementary kind can maintain itself, as
long as the simple conditions of life, to which it is adapted,
are not changed. Thus some of the lowest
animal forms, such as infusoria and rhizopods, seem to have existed
for infinitely long times almost as they now are.
The adaptation to conditions is here as good as complete.
The same may be the case with beings
that at an earlier stage of development have possessed more
developed forms and organs and therefore probably
also have had a richer life of feeling. Animals that in their
youth live freely, but later lead a parasitic existence,
are wholly absorbed in the life of nutrition and lose all nerves and
muscles that do not directly serve this. --- Also among
human beings there are examples enough of an adaptation to
conditions that hinders further development. Of the Fuegians
Darwin remarks in his description of their (in our
eyes) sad existence\footnote[1]{Rejse om Jorden. p. 242.}: ``There is no reason to
believe that the Fuegians are diminishing in number; therefore we must
assume that they enjoy such a measure of happiness, however
% --- p. 33 ---
\setcounter{footnote}{0}%
\opage{33}it may be constituted, that life acquires value for them.
Nature, which made habit an irresistible power
and let its effects become hereditary, has fitted the Fuegian
to the climate and products of his wretched country.''
Also higher-standing primitive peoples, and not only
primitive peoples, present similar examples. It is need
that drives man onto the path of progress. Where
even wretched conditions have not made adaptation impossible,
the life of feeling is conditioned only by a narrow circle of conditions.
There can then perhaps be attained a more complete contentment
than under more manifold and more complicated conditions. A
small vessel is easier to fill than a large one\footnote[1]{Vas minus majore plenius esse potest, qvamvis liqvoris minus contineat. \emph{Cartesius}: Epistola IV. (De vita beata).}.

It might now seem that it must be an essential
objection to the principle of welfare, that it must lead
to recognizing such an existence, confined to the first necessities of life
and to purely elementary feelings, as standing just as high
as the existence that is filled
with mental work and cultural activity, --- indeed, as
standing higher than this latter, which will hardly bring
so unmixed and secure a well-being, but be connected
with ever new want, new striving and new unrest.
--- An existence like that of the Fuegians probably strikes everyone as
miserable and wretched, because they surely often suffer privation and
want. But we can also take another example.
Alexander von Humboldt met in South America a population
that lived on banana trees, which are so fruitful that
a tønde of land planted with them can feed
50 individuals. These trees require no real cultivation;
only the ground around the roots must be worked with a
hoe once or twice a year. These people
can, in the double meaning of the word, ``nourish'' themselves with very
% --- p. 34 ---
\setcounter{footnote}{0}%
\opage{34}little work. Therefore they are also dull and without culture.
But the needs they have, they get satisfied.

What for us would make such a life unendurable,
\emph{the urge to activity, development and}
\emph{progress}, does not exist at such stages. This
urge is, as must well be noted, \emph{itself a consequence
of development and progress}. --- One of Darwin's
predecessors, \emph{Lamarck}, assumed an inner, inherent urge to development
in all living beings. But \emph{Darwin} has,
relying on experience, maintained that development is always initiated
by the working of external conditions. For Lamarck it became
a difficulty how the very simplest forms
have continued to exist, --- why by virtue of the
assumed urge to development they have not long since passed over
to higher forms. He solved this riddle by assuming
a generatio æqvivoca, a new arising of the simplest
organic beings out of inorganic matter, an assumption
which has so far not been confirmed. For Darwin the continued existence of
the simplest forms causes no such
difficulty. For according to his doctrine only those
variations are favored that can be of benefit in the struggle for
existence under the given conditions; qualitative selection therefore does not
necessarily entail a progressive development\footnote[1]{Arternes Oprindelse. p. 153 f.}.
As regards human beings \emph{Th. Waitz}
has clearly shown how the urge to development is itself a product of development:
``Of a people without intercourse and without
rivalry with others, a people which with relative ease
or only with overcoming of accustomed difficulties
regarded as necessary can satisfy its
need by what immediately surrounds it, and
which feels content therewith and lives happily, --- of such a
people one must not expect that it shall make any
% --- p. 35 ---
\setcounter{footnote}{0}%
\opage{35}effort whatever to civilize itself. He who
has what he needs and feels well thereby in every
respect, he does not work. One does not become civilized
without necessity and of one's own accord, or perhaps in order to follow a
noble urge that dwells in the heart. Is it otherwise then
in our modern society? Does there not first belong a long
school of cultivation, to instil in man interest
in work as such? Let one reckon up
among the so-called educated, how many make
an effort for their own or others' education's sake
without being compelled to it! One will then cease to
doubt the equality of the savages with us, as far as mental
faculties are concerned, at least for the reason that they omit
to exert themselves without necessity.''\footnote[1]{Die Indianer Nordamerikas. Leipzig 1865. p. 69.}

Peculiar to the state of nature in contrast to
the state of culture (in so far as one can sharply set these
two states in contrast to each other) is that over there
the impulse to change in mode of living and way of thinking
must always come from outside, whereas here --- even if it
should be in comparatively small circles --- a
stronger or weaker urge to advancement stirs, which is always
kept up. This difference forms an analogy to
the difference between inorganic and organic existence. In
the organism the play of forces is maintained with a certain
independence of the influences of the moment and of the limited surroundings;
in the inorganic world (and likewise
when one compares the lower organisms with
the higher) external influences play a far more decisive
role. Among cultured peoples there has been awakened an urge
to change and differentiate, to shape life in all
directions and to sharpen it within each single
direction. Mental feelers have grown up which are in
% --- p. 36 ---
\setcounter{footnote}{0}%
\opage{36}uninterrupted activity. A new kind of feelings now becomes
possible, feelings whose nature is determined not only by the
ends that are attained, but which also attach themselves immediately
to the work and the activity itself. Man has thereby
become more independent and freer, and his mental life has
gained in depth and strength, in that it is no longer
determined merely by the external world, but essentially by the
forces that stir in the inner world. Higher
demands are now also made on life.

\emph{The arising of this urge to development can be quite
in accord with the principle of welfare}. For that stability
in the state of nature, which may strike us now as
pitiable, now as paradisiacal, is dependent on the stability of external conditions.
But an absolutely assured stability is not given in
nature. Even if the nearest surroundings are not changed,
changes occur elsewhere in the world and among other
beings, and the struggle for existence can then force the
stable beings either to perish or to develop
the capacity for progress in themselves. The more highly developed beings
often in the struggle for existence put the lower developed
beings to this choice. It is put to many primitive peoples, or
cultured peoples, that have remained standing at a lower stage;
they so often go to meet their doom when the higher
culture approaches.

What has here been shown by examples from races and
peoples can also be shown for the single individuals.
A perfect adaptation to limited
conditions always includes a danger: the danger that
the individual, when its conditions of life and its horizon are widened,
shall not have the inner conditions for asserting
itself. Childlike naivety, dreaming fancy,
sensuous enjoyment each have their right and their significance;
but the state easily becomes a sleepwalker's state, in
which there is only security and happiness until the
% --- p. 37 ---
\setcounter{footnote}{0}%
\opage{37}awakening occurs, and where the greatest helplessness can set in
at the awakening. And here we leave entirely out of consideration
that such states often exist only at the expense of other
beings, in that these either must give the sleepwalker
help or at least do without his slumbering forces.

Welfare can therefore not consist in a passive state, produced once
for all, a state that is not itself
the point of departure for new development and progress. Welfare
must in the highest sense consist in \emph{a state}, in which
force has been gathered up and rich possibilities laid down for the future,
and which engenders an urge to set new ends and
to initiate new striving. It must consist in a state,
\emph{that both in and for itself is desirable and engenders
striving after new desirable states,
thus a state that is not only end,
but also means, not only has value as
effect, but also as cause}. The feeling of pleasure
attaches itself here immediately to the activity and
the work itself, to the unfolding itself of the forces, not merely to
the result. Where such a feeling of pleasure is possible, there
many sufferings become endurable, which on a lower
level would be signs of the dissolution of life itself.
Tension and longing, disappointment and privation cannot
fail to appear and will with a certain rhythm succeed the cheerful
advance toward the goal one has set oneself; but through
all oscillations the basic direction and the basic impulse
can be preserved.

This is the ideal that the principle of welfare, when it is thought through,
calls forth for us. How far it can be attained is
a question which at each given time and for each single
individual must be answered by an experimental way. Very often what is at stake is
not only changes in the nature of the individuals,
but also in the ordering of society or even in the external
nature, if an approach to the goal is to take place.
% --- p. 38 ---
\setcounter{footnote}{0}%
\opage{38}But for us in this context the main point is that we
here have a standard for the assessment of actions and
institutions,

This standard corresponds to a tendency that runs
through the whole of organic nature, where pleasure on the whole
and in the large indicates life and advance, pain decline
and death. The principle of welfare maintains \emph{the right of life}:
every being has a right to exist and develop, in
so far as thereby no so much the greater pain is produced
for other beings. The ethics that builds on
the principle of welfare thus seeks to continue the development of nature
in a more conscious and harmonious manner. It calls
for seeking means to life and advance which
the unconscious development of nature would not lead to, and it
strives to exclude the unnecessary pain and distress
which the blind struggle for existence brings with it. It
comprises a series of tasks extending all the way from
the alleviation and help which compassion motivates, to
the work for the realization of the highest social, aesthetic and intellectual
ideals. It is the task of special ethics
to investigate how these various tasks
are to be treated, and how they relate to one another.

5. But because welfare consists in activity and development,
one cannot conversely say that activity and
development are necessarily connected with welfare or
lead to welfare. Because it is incorrect to believe that limitation
of the needs always leads to the goal, it by no means
follows that unlimited multiplication of the needs
should be the right thing. Cultural development
can assume forms and strike out on paths that do not agree
with the principle of welfare. We therefore also find in history
at times sharp and penetrating warnings against the
existing culture. Thus in Greece by Socrates,
% --- p. 39 ---
\setcounter{footnote}{0}%
\opage{39}the Cynics and the Stoics, in the eighteenth century by Rousseau
and in our day by Leo Tolstoy. The opposition of such great
minds must in every case make us attentive.


I here set aside the question whether that which we call
culture can without more be imparted to any people. Even
if one has been too quick to deny primitive peoples
the capacity to receive civilization, it does not follow that
it would be a good for a people to give up the forms of life developed in and
through itself, in order to let itself be regulated
according to forms and ideals that are developed under wholly
different conditions. So directly cannot even the most justified
and perfect culture be imparted. \emph{Waitz}, who expressly
maintains that no proof has been given of the Indians'
incapacity for culture, nevertheless praises those Indian chiefs
who opposed the penetration of European civilization
into their people, for ``love of their people
and right insight into its true welfare.'' There lies here
(as in the clash of ethical systems) a possibility
of tragic conflicts, of which the history of the Indians in particular offers many examples.

The reason why conflict can arise between culture
and welfare lies in the restlessness and recklessness of the striving for culture.
It is with it as with the involuntary
(``spontaneous'') urge to movement, which leads to
using forces and limbs, merely because energy is present,
without regard to a valuable end guiding it,
so that it becomes accidental what result the movement
gets. The striving for culture can lead partly to overstrain
of the forces, partly to one-sided application of them
(owing to excessive differentiation), partly to isolation,
to a fragmentary separating out of the individual activities.
In the single individual single faculties are cultivated and developed
(in some, e.g., the intelligence, in others the physical
% --- p. 40 ---
\setcounter{footnote}{0}%
\opage{40}capacity for work) at the expense of others; harmony, the capacity
to feel oneself as whole and unity, is lacking. Such
one-sidedness makes the individual acquire value only as a
wheel in the great machinery, become merely a means, not an end.
It is connected with a one-sided social development;
the suppression of single sides in the nature of the individual
is caused by the suppression of single estates
and classes in society. If one wished to identify
culture and ethics and to regard development as a surer
criterion than welfare, one would wholly overlook that there
is a social question. The social question is
an ethical question and at the same time a question of the justification of our
culture, both by virtue of the principle of welfare.
Would it be right that the products of material
and ideal culture should benefit only a small
minority, while all others are excluded from
sharing in them? That would quite conflict with the ideal
of a society that follows the principle of welfare. For the
greatest welfare is present where each single individual
develops in an independent manner such that it thereby
in the best way helps other individuals to a similar
development from their peculiar stage. There exists
then \emph{a harmonious society of independent personalities}.
This is now the highest ethical idea. Each
single individual stands as a little world for itself, and yet
in the most intimate reciprocal action with the great world,
of which it is a part. The individual serves the race,
and the race serves the individual. Every special position, every
inequality in distribution must be justified by what the
various conditions of life and tasks require, and the faculties
and impulses of every individual are to be developed as fully and richly as
is compatible with what the life of the race as a whole requires\footnote[1]{See further on this ``Etik''. p. 31. 164. 227---240.}.

% --- p. 41 ---
\setcounter{footnote}{0}%
\opage{41}6. It follows from what has been developed here that the principle of welfare
by no means is so easy to apply in the single
cases as it may perhaps seem. The special
connection can be so complicated that we cannot
survey it or foresee the effects of our intervention.
We cannot deduce from the principle of welfare a priori a
system of single actions or an ordering of society
or an entire culture. Its significance is (like that of
the principle of causality in the theoretical domain) to pose
problems and to guide in their treatment. It is
regulative, not constructive. It presupposes the
immediate, involuntary life in the individual and in society
and steps in only when there takes place a conscious
discussion of the value of what has thus been developed and of the manner
in which it is to be carried further. All ethics thereby acquires a
historical character. For we never begin --- neither as regards
the single individuality nor society --- absolutely
from the start, but have to work further on a given basis
under the guidance of the principles and ideals that spring
from our nature.

% --- p. 42 ---
\setcounter{footnote}{0}%
\essay[The Law of Relativity in Ethics]{The Law of Relativity in Ethics.}
\opage{42}


\begin{center}I.\end{center}


The aim of this treatise is to show that the validity
of all ethical laws rests on definite relations and conditions,
so that the law of relativity, the fundamental law of our knowledge,
finds application in the domain of ethics too. It
will appear that this view serves in essential respects
to characterize the nature of ethical knowledge and of the
ethical laws. ---

We may begin by examining in what sense and with
what right one speaks of ethical laws at all. The word
law is not used here in the same sense as in natural
science, in psychology and in sociology. For ethics does not seek
to demonstrate the rules which human willing and acting actually follow,
but the norms and ideals which it \emph{ought} to follow, the standard
to which it ought to be subject. The sense of the word
law that underlies philosophical ethics therefore recalls
the way in which the word is used in positive jurisprudence
and also in theological ethics,
as well as in the science of positive morality
(the norms lying in custom and usage and in existing public opinion).
But from these sciences philosophical ethics differs
in that it builds on no
% --- p. 43 ---
\setcounter{footnote}{0}%
\opage{43}external, natural or supernatural authority, but seeks to
ground its norms and ideals on human nature itself,
as we know it from experience.
It has been thought that this circumstance made it inadmissible
to speak of ethical laws.

According to \emph{John Austin} a law presupposes two
parties, one who gives the law, and one for whom the law
holds, and the binding force of the law lies in the fact that the one
who gives the law has the power to maintain it by
attaching pain to its transgression. Since the law thus
maintained by power produces uniformity in the manner of action,
the individual actions being determined
by the demand expressed in the law, the word law can at the same time
come to mean a general rule, which is followed in
the individual cases, and it is this side of the concept of
law, and only this, that we make use of when we speak
of laws of nature. We mean thereby the regular sequence
of the phenomena of nature, without our needing to represent
to ourselves this regularity as the effect of the power of a commanding
will. Even if this usage can be
permitted, Austin flatly denies that the concept of law can
find application to the rules of action which the
individual human being lays down and acknowledges for himself.
For here there is neither a law of nature in question, nor
is there a power standing above the individual which
could maintain the law in cases of transgression. ``For though he
can form the sincere resolve to inflict pain on himself,
if his conduct departs from the guide he has laid down
for it, the actual infliction of this possible pain depends on his own will.''\footnote[1]{Lectures ou Jurisprudence. I. p. 213 f.}.

Might there not, however, be a \emph{psychological
basis for the concept of ethical law}, that could
% --- p. 44 ---
\setcounter{footnote}{0}%
\opage{44}contain the three elements which according to Austin are presupposed
in the strict application of the concept of law: 1) a relation
between superior and subordinate, 2) a general
rule of action and 3) the maintenance of this general
rule by pain in cases of transgression?

Various feelings stir in man,
each determined by its content of sensation or of idea.
There is a feeling which is bound up with the individual's
own self-preservation. And there are other feelings which let
the individual share in the weal and woe of others. Since now every
feeling is at the same time a motive of action or engenders an impulse
to action, there can appear in man's consciousness an impulse
to action which is determined by
regard for a more comprehensive order of things, for the
welfare of a larger group, over against an impulse to action
which is determined only by regard for the individual's own
limited interest. There will then be something in the individual
which points beyond the individual itself and assigns it a place as
a single member of a great whole, so that the welfare of the latter
may at times demand the sacrifice of what regard for
the individual's own welfare as an isolated being would require.
Here, then, there will be in the individual's own nature an opposition
between two parties, between an impulse which is bound up
with the idea of a greater order of things, and an
impulse which has a narrower horizon. The first element
in the concept of law is hereby demonstrated. --- In reality
matters stand no differently, even where the party
that gives the law is an external authority. For this authority
must after all make itself known to the individual by awakening
its feeling (whether this feeling be fear or
reverence), and the law will here too reveal itself
through a relation between two feelings or impulses in
man's interior. For one who can eradicate the feeling
of fear or reverence in himself, the law does not exist. It
% --- p. 45 ---
\setcounter{footnote}{0}%
\opage{45}can be seen from this that philosophical ethics has just as much right
to apply the concept of law as jurisprudence, theology
and positive morals. --- I do not go further here
into how this inner opposition between a
more comprehensive and a more limited regard within
the life of feeling and of impulse comes about. It is an opposition
which has developed historically under the influence
of social life on man by virtue of general
psychological laws.

Thus the first of the three elements in the concept of law
can very well be found again in the laws which philosophical
ethics sets up. The same now holds also of the one of
the three elements which Austin could least of all find
here, namely the sanction, or the suffering brought about by transgression of
the law. Austin holds that it depends merely on the individual's will
whether it comes to suffer pain
by transgressing the law that was supposed to spring from
its own nature. But when the ethical law is the formulation
of an impulse which is called forth by the individual's
feeling being determined by regard for a larger whole, then
such an impulse cannot meet resistance and contradiction
from other sides of the individual's nature without more or less intense
pain arising. When the impulse
which causes the individual not to obey the ethical law it
has itself acknowledged is unable fully to suppress or eradicate
the impulse which led to acknowledgment of the law, then
there will arise the inner disharmony, the painful feeling
of inner dissolution and degradation, which we call bad
conscience or also repentance. And it does not depend on
the individual's arbitrary choice whether this pain is to set in
or not; it follows with psychological necessity
from the conditions at hand. What expresses
itself in bad conscience and in repentance is not anything
accidental or arbitrary. It is, under the given
% --- p. 46 ---
\setcounter{footnote}{0}%
\opage{46}circumstances, an unavoidable psychological result --- just as
unavoidable as external punishment is, when the power of the state
is strong, wise and swift enough to lay hold of the transgressor
of the civil law. But the power of the state is, after all, not always able
to strike the criminal with its punishment, and
so it can likewise happen that repentance fails to appear, namely
when the impulse which leads to the acknowledgment of a
larger whole is extinguished, or when it is not strong and
vivid enough to assert itself in opposition to
the impulse which led to the performance of the act that was contrary
to that acknowledgment.

Finally, the property of the law that it
denotes a general rule can also be demonstrated in the ethical
law. In the idea of the larger whole to which
man feels himself bound --- the family, the clan, the nation,
humanity --- he has a principle from which
general rules for his willing and acting can be derived. If the individual
is not to contradict himself,
he must judge the state and direction of his will from
the most universal and exalted standpoint he knows,
from the thought of the great whole of which he feels
himself a single member. The particular rules follow with logical
consistency from this thought, once it
has been acknowledged. Ethical science is then only an
unfolding of what is implicitly given in universal
sympathy, as this has developed historically
in the human race. ---

But when it has thus proved justified to apply
the concept of law in ethics, the preceding
argument also implies that both the general
principle and the rules which are derived from it hold only
\emph{in relation to} and \emph{on condition of} a definite psychological
basis. If there should be human beings
who lack this basis, they neither could nor would
% --- p. 47 ---
\setcounter{footnote}{0}%
\opage{47}acknowledge the ethical laws. How one must take up a position
toward such individuals, if they exist, is a matter by
itself, which we shall come to touch on later; here we assert
only that all ethical laws hold only in relation to a
definite psychological presupposition. This is the \emph{psychological}
relativity of ethics\footnote[1]{Cf. my ``Etik''. p. 35---37. --- See above p. 7---17.}.

\begin{center}II.\end{center}


But it is not sufficient to establish the psychological
basis and to set up the corresponding ethical
principle. The relations under which the rules following from
this principle are to be applied are so intricate and
complex that it is difficult, indeed impossible, to find
rules which could without further ado be applied
in all \emph{special relations and situations}.

The \emph{general} law or rule cannot foresee
all possible situations in which a human being may be
compelled to take a decisive resolution. It is a
difficulty which the ethical law shares with the
juridical and the theological law, as well as with the law of
positive morality. Life is too rich and manifold
to be put into a fully finished system. The relations
pass over into one another through many different nuances and
intermediate stages; sharply drawn boundaries do not exist.
The general law, which is to fit \emph{all} cases
of a certain kind, easily becomes inapplicable in the individual
determinate case, which may be altogether unique.
The ethical law must therefore --- if it is not, for the sake of an apparent
formal consistency, to sacrifice the real
right of life --- fix the general \emph{tendency} or \emph{direction}
which the will is to follow, and not give special
% --- p. 48 ---
\setcounter{footnote}{0}%
\opage{48}instructions as to how in every individual case one is to act unconditionally.
Only in this way can the universal validity of the law be reconciled with
the relations, which differ at different times and
in different places. This one may call \emph{the historical relativity} of the
content of the ethical laws\footnote[1]{See my ``Etik''. p. 160 f.}. Proceeding from the same psychological
basis and in accordance with the same general principle of valuation there will
under different historical relations, in different situations,
be different decisions to make with regard
to the ordering of society and the conduct of the individual.
What is possible at one stage of development is not
possible at another. From the general
principle of welfare (any more than from any other principle)
no ordering of the state or the family can be constructed which
could without further ado be introduced everywhere. Every social order
presupposes that those who are to be the members of society
live under certain definite inner and outer conditions.
On bare ground no building can be erected.

The ethical laws do not hereby lose their significance.
They cannot tell us how on the sea of life we are to
manage each single wave that rises before us,
or how we are to tack in the individual cases.
But they can tell us in what direction we are on the whole
to steer, and they make it possible for our compass to
guide us without deviation, even if the road be ever so
winding and irregular. They formulate the general
demands that the highest end we acknowledge makes
of us. The means of satisfying these demands may
shift within wide limits. But when the general
direction of the will is given, and when there is sufficient
strength and consistency in character, the individual
ways and means too will be found, without one's coming into
% --- p. 49 ---
\setcounter{footnote}{0}%
\opage{49}conflict with the principal end. \emph{Kant} has rightly said
that casuistry is not a science but an art, ``not
a doctrine of how something is to be found, but practice in
how the truth is to be sought''.

Casuistry concerns the problems which the intricate nature of the given objective
relations poses, when a decision in accord with the ethical principles is to be
made. It seeks to analyze the relations into their elements in order
to find out what the consequence of the general principles
must be under these particular relations. It sets up
general rules for cases of conflict or seeks
by means of invented examples to exercise the ability to find what is
right. But analysis and general rule cannot always
be transferred from one case to another. It
may e.g. be correct that a higher, more comprehensive
end is to be fulfilled in preference to a lower, more limited one;
but the difficulty is precisely to decide what in
\emph{this particular} case is the most significant
end. Setting aside the varying multiplicity of the cases,
all ends hang together in such a way that one cannot
intervene in one of them without its having effects on
a great number of the others; and how then to cast up
with certainty the account of what will prove to have the preponderance on the whole, the beneficial or the
pernicious effects of a certain kind of intervention?

One cannot say without further ado that the more comprehensive
regard in the individual case is to have the preponderance
over the limited regard, so that e.g. where
the interests of the state and of the family come into conflict, the family
ought always to yield, or my interest is always to be subordinated,
where it comes into conflict with those of several others. For the
sound development of life depends not only on its extent,
but also on its strength and inwardness, not only on extensity,
but also on intensity. In the narrower
% --- p. 50 ---
\setcounter{footnote}{0}%
\opage{50}society a more inward and fuller life can be lived
than in the large society, and the very force that holds
the large society together unfolds and is exercised in the smaller
societies, so that the vital strength of the former depends on that of the latter. Within
the narrower society there arises a fellow-feeling, which after
having grown strong can be extended to become an
essential part of what supports the larger societies. The
real productivity takes place in the smaller circles;
from them the new thoughts and the new forces proceed. Therefore
regard for extent cannot without further ado become
decisive in cases of conflict. Something similar holds,
as regards the individual's self-preservation. The life of the
race gains in strength the more and the better the individuals
are able to preserve and assert themselves in
material as well as in mental respects. Within individual
ethics, therefore, self-assertion is not without further ado
to be sacrificed to devotion, any more than within social
ethics the interest of the narrower society is always to be
sacrificed to that of the larger. --- But how difficult may it then
not be in the individual case to make the right decision,
so that neither the strength and inwardness of life nor
its extent and diffusion suffers harm? ---

Another case in which there may be use for
casuistry is the relation to other people's prejudices.
Here Bentham\footnote[1]{Of the influence of time and place in matter of legislation. (Works. I.). p. 180.} has given the following rule: ``He who
frivolously and without necessity attacks a prejudice, and
he who submits to being led like a blind slave by
prejudice, both go astray from the right path of reason''. ---
But how can one decide in the individual case
whether there is necessity present for an attack, or
whether it would be unworthy slavery to accommodate oneself to
% --- p. 51 ---
\setcounter{footnote}{0}%
\opage{51}the prevailing prejudices? Here the relations vary in a
particularly high degree, since the question is whether the
individuals we face have the conditions for fighting through
the (perhaps both outer and inner) struggle into
which we plunge them by attacking what has hitherto
stood immovably fast for them, and what perhaps
has given their life support and strength. ---

Bentham thinks, to be sure, that in our days one is too inclined
to deny that the laws which for our time are
the best possible would also have been so in earlier
times. There are laws which would have been beneficial
at all times. But he expresses the principle of historical relativity
in a striking manner in the following words\footnote[1]{Op. cit. p. 193.}:
``If I were to choose what I would with greatest truth
and readiness ascribe the magnificent qualities
of generality and immutability to, it should rather
be certain \emph{grounds} of the laws than the laws themselves: rather
the \emph{principles} on which the laws are to be built,
the subordinate \emph{considerations} which follow from these
principles, the best plan according to which they can be
combined, the \emph{regards} by which the legislator ought to let
himself be guided, --- than any particular laws whatsoever which
the legislator ought to give, in order that those who are entrusted to
his care might be well governed''.

\begin{center}III.\end{center}


1. Besides \emph{the psychological basis} and
\emph{the historically given situation}, there is yet a
third factor to be taken into account, a third relation, by
which the ethical decision in the individual cases must be
determined. It is this third point which it was in particular
% --- p. 52 ---
\setcounter{footnote}{0}%
\opage{52}my intention in this treatise to call attention
to, because it has received less notice than
the first two (especially than the second). It arises from the fact
that different individuals on the same basis (and therefore
with the same ethical principles) and in the same situations
meet with \emph{different predispositions and capacities}. Just as
little as the same principles lead to the same consequences
in different situations, just as little can they in the same
situation lead to the same consequences, when the willing and
acting individuals are differently endowed. We stand
here before \emph{the individual relativity} or the
personal equation in ethics\footnote[1]{In my ``Etik'' p. 133---134 I have already emphasized this point. It is expressed more definitely in the German edition (Leipzig 1888). p. 154 f., where I have designated the law of individual relativity as the law of relativity in ethics. The concept of a law of relativity may, however, be given a wider sense, so that the psychological and the historical relativity also come under it.}.

The different ethical conceptions seem on the whole to
agree that the ethical laws or commands are the
same for all in the same situations. The ethics which lays
the principle of authority at its basis proceeds from the assumption that the will of authority
demands the same of all individuals. On the
one side stands the eternal will, on the other side
men, who are to obey its commands. In a similar
way the laws of the state and the commands which custom and usage impress,
are the same for all individuals. There holds, as
one says, no respect of persons. According to the morality of reason
a general law arises from human
reason, which everyone must follow at the risk of coming
into conflict with himself, that is to say: with the universally human
in his nature. The adherents of the principle of welfare too
have as a rule been inclined to conceive the matter in
a similar way. The ethical laws they derive by deduction
% --- p. 53 ---
\setcounter{footnote}{0}%
\opage{53}from the principle of welfare, combined with the given
situation, and these laws are then to sound alike for all.

That there must be an error in this assumption can
be seen in the following way. \emph{On account of the individuals'
different capacities and impulses, one and the same demand will take on a highly different
significance} for them. The demand for a certain degree of self-control
with regard to some natural impulse will thus be able to mean
something highly different for the different individuals,
in that each of these may have the natural impulse constituted in a different
degree. What the one can do without the least
effort, because there is an immediate harmony between
the demand and his original nature, the second can do
perhaps only after a serious, wearing struggle, whereas
the third perhaps is altogether unable to overcome
his unruly nature. What for the first
accords with the natural inclination itself becomes for the
second an act of duty performed with strong exertion and conscious willing,
and signifies for the third a task which
he cannot solve, since from the first he lacks
the conditions for it. The law is thus in \emph{reality}
not the same for all three. And now it is of no
use to distinguish between the law in itself and the individuals'
different capacity to fulfill it. The ethical law differs
after all from a theological or juridical law in that it
springs from human nature itself. It does not come
from without, but is to unfold for the individual out of the
innermost of its being. It cannot then be attained merely by
deduction from principles and situations which are common,
and which contain no regard for the peculiarities
of the single individual. So long as only those two points
are included in the deduction of the law, it cannot be the
individual's own law, but an abstract impersonal demand.
\emph{The truly ethical law must not only}
% --- p. 54 ---
\setcounter{footnote}{0}%
\opage{54}\emph{be directed to the individual, but also be individualized
in such a way that the individual's being comes
to development through satisfaction of the
demand of the law}. It must take regard of the special
starting points for willing and acting, which lie
in the nature of the single individual. It would otherwise demand
something different, and more or less, of the one individual than of
the other: only when it sounds different for each single
individual according to its possibilities does it in reality make
the same demand of them.

This follows quite simply from the principle of welfare. The
single individual is not merely subject to the law which declares
what the welfare of the race demands, but is itself a peculiar
member of the race. In the law which holds for
it, therefore, two regards must meet in a harmonious way:
the claims of the race and of the individual. The law has not been found, is
not the right law, so long as not both these regards
have come into their full right. In striving and working for the welfare of the race,
the individual must
at the same time work for its own welfare as part of
the race. Therefore he must have his special ethics, which
both qualitatively and quantitatively is different from that of
the other individuals, in that his nature and dispositions become
co-determining of the way in which the law
shapes itself for him.

It would be a great misunderstanding if one
were to draw from this the conclusion that all self-sacrifice was
to be rejected. The self-sacrifice which springs from enthusiastic
love, from free resolve, is precisely an expression of the individual's
own innermost being, an outlet for a deep urge
in it. In ideal self-sacrifice self-assertion and
devotion go in one: life acquires, strange as it sounds, only
value by being sacrificed. (See p. 23. 29ff.). The question is precisely
in the individual cases whether there is capacity in the individual for
% --- p. 55 ---
\setcounter{footnote}{0}%
\opage{55}such self-sacrifice. This kind of action can no more
than other kinds of actions be demanded with equal
ease of everyone. Here the problem emerges, since the capacity
and the impulse to enthusiasm and devotion are in fact not
the same in all.

Just as, according to Fechner's law, sensation does not
depend on the absolute strength of the impression, but on the relation
between the present impression and the preceding one,
which has already affected the sensing individual,
so from a purely ethical standpoint the demand
which is in reality made on the individual must not be one
with the demand which --- without regard to the individual's
presuppositions --- could be derived from the general principles
and the given situation, but depend on the relation
between this purely objective demand and the individual's presuppositions
or dispositions, and this relation will be
different for the different individuals. The purely objective
demand may be compared to a burden, which in itself
has the same weight, but which is felt differently on
the shoulders of different individuals.

It is thus a complete individualization of the
ethical law which is hereby demanded. The formulation which
can fit all, or at any rate a great number of individuals,
can only be coarse and superficial. In the determination of
$\pi$ we can at times content ourselves with $\frac{22}{7}$, whereas we
by making greater demands on accuracy arrive at
several hundred decimals, and can still always think of
greater accuracy. Just as in the psychological
description and understanding of a character there is always something
that escapes us, because it is impossible to grasp all
nuances and bring out the individual qualities in their
whole multiplicity, so too in the determination
of the individual task there will be individual
% --- p. 56 ---
\setcounter{footnote}{0}%
\opage{56}nuances for each single human being, which cannot be
represented exhaustively in any formula. Or rather: into
the same formula there must be inserted, for each single individual's
part, particular coefficients. The juridical
laws cut everyone to one pattern, and the current
ethical laws do so too to far too high a degree. A fine
sense for the individual differences in qualitative and
quantitative respects is still extraordinarily rare.

The same line of thought which led ethical valuation
to go back from the action to the motive must consistently
be carried a step further and take regard of the predispositions and
dispositions which condition the arising of the motives. The possibility
of a motive's arising is not the same in
all individuals. Only an absolutely indeterministic conception
could suppose that (as a German jurist\footnote[1]{Quoted in Vierteljahrsschrift für wissenschaftliche Philosophie. VI. p. 201.} has put
it) ``the moral capacity to inhibit lawless impulses
is an unchangeable quantity, which is normally to be posited
in all human beings without regard to individual predispositions''.
Such an assumption conflicts to such a degree with all
psychological experience that it will surely become rarer
and rarer in scientific circles. --- And besides,
it is not enough that a motive can assert itself
with a certain strength. The influence it can gain on
the will will depend on the strength of the other motives in relation
to it. Thus, as Butler\footnote[2]{Sermon, XII.} already remarked,
two persons may have exactly the same degree of
compassion; but in the one there stirs a very
strong ambition or desire for revenge, in the other
these feelings are very weak; they will then be placed
highly differently in ethical respect. It is here as
% --- p. 57 ---
\setcounter{footnote}{0}%
\opage{57}with a pair of scales: which scale is to sink depends
not on the weights that are placed on the individual scale,
but on their relation to the weights that are placed on
the other scale.

2. The first thinker who has directed his attention
to the ethical problem here brought forward
is \emph{Aristotle}. In presenting his well-known doctrine
of virtue as the right mean he expressly impresses
that objectively considered (κατ’ αὐτὸ τὸ πρᾶγμα) the mean is
the same for all, in that two extremes can be indicated,
from both of which it is a matter of keeping equally far away,
but subjectively considered (πρὸς ἡμᾶς) it stands differently on account
of the individual dispositions. Just as little
as one and the same quantity of food suits the practiced
athlete and the first-time beginner in gymnastics, just
as little will for all the inner midpoint of the life of the impulses
be situated at the same place (τοῦτο δ’οὐχ ἓν οὐδὲ
ταὐτὸν πᾶσιν)\footnote[1]{Eth. Nic. II, 5. p. 1106 a, 32.}. What in the one must be called
courage, in that in him it presupposes the subduing of
a strong timidity, that can be practiced altogether without
exertion of will and self-conquest by the other.
This way of regarding things Aristotle carries through in an
interesting manner with regard to the individual virtues. ---
This Aristotelian doctrine, which testifies to a surprisingly
fine eye for the significance and right of the individual personality
in ethical respect, stands, as I have pointed out in another place\footnote[2]{``Etik'', p. 133 f.},
in a strange disharmony with
the altogether predominant place that Aristotle gives
to the state in relation to the individual. Aristotle has made
no attempt to show how that apparently
ultra-individual doctrine can be reconciled with the objective,
% --- p. 58 ---
\setcounter{footnote}{0}%
\opage{58}social way of regarding things. Socially there is demanded --- so
it seems at least --- often with necessity more
of the single human being than he can yield according to the ``right mean''
determined by \emph{his} nature. In the relation between
individual and society there emerges a series of problems
for which the ancient thinkers as yet had no
clear eye. Nevertheless that Aristotelian doctrine stands as
one of the most ingenious points of view that has ever been
advanced in ethics\footnote[1]{In more recent times the thought of the individualization of the ethical law has been advanced by \emph{Jakobi}, \emph{Schleiermacher} and especially by \emph{Beneke} (Physik der Sitten. 1822), and in our own literature by A. S. \emph{Ørsted} in his beautiful treatise ``om Grænserne mellem Teori og Praxis i Sædelæren'' (1811), which is reprinted in the first volume of ``Eunomia''.}.

One might perhaps think that it is a problem insoluble in and of itself
which Aristotle, without knowing it, has
posed by his doctrine of the individually determined ``right
mean''. It might seem that a choice must be made
between two equally unfortunate ways out. The one way out would be
an ethical subjectivism, according to which every generally valid
ethical law falls away, since the law becomes altogether
different for each single individual. The other way out
would be an irreconcilable conflict between the single human being's
individual ethics and the ethical
claims of the social life of man. --- The conception of the ethical law which
has been advanced in the foregoing nonetheless helps us out of
this dilemma.

Subjectivism is repelled by the fact that, despite the different
forms the ethical law must assume in and for the different
individuals, the \emph{direction} of their striving
can nonetheless be the same. The ethical laws formulate, as we have already
seen from another standpoint, tendencies or
% --- p. 59 ---
\setcounter{footnote}{0}%
\opage{59}directions for the conduct and ordering of life. When several
ships tack up against the wind, they attain --- on account
of their different build --- with equal exertion
different results; but they can nonetheless all keep
the direction; this lies in the continuation of the course they
hold. The objective way of regarding things need not
therefore fall away because one maintains the individualization of the ethical law.
The individual differences in predispositions
and impulses ethics must take as factual starting points,
which cannot be erased, but from which it is a matter of
bringing out the most possible. They entail, to be sure,
a barrier for ethics as a science, but a barrier,
which bears witness not merely to the limitation of our thinking,
but first and foremost to the rich multiplicity of life,
which cannot be fitted under definite rubrics, although
what stirs in the many individual existences
can nevertheless be directed toward the highest goal of all human
acting. This goal is at the same time so exalted that
the significance of the individual differences for the ideal
way of regarding things is restricted to a high degree, if it does not
altogether vanish.

The other apparent consequence falls away,
when one thinks through more closely the relation between individual
and society, and when one at the same time grasps the
real significance of the ethical laws in the right manner. ---
The single individual is from first to last a member of society,
and \emph{the life of society is the life that stirs in all the
single individuals}. Society's goal is therefore not
reached so long as its demands on the single individual
bring about an irremediable disharmony in the latter's interior.
When the will of society cannot be carried through without the
single individual's being cowed or broken, that points to
an imperfection in the very order of society. The ideal
is reached only when the work the individual performs in
% --- p. 60 ---
\setcounter{footnote}{0}%
\opage{60}society's service at the same time serves for the unfolding of its own
capacities and impulses in a free and harmonious manner. The
greatest welfare is present only where every single individual
develops in an independent way and precisely thereby
in the best manner helps forward a similar development in
the other individuals. Thereby he comes at one and the same time
to stand as goal and means. What is decisive is that
in the interior of the single individual a work is performed in the \emph{direction}
of the ethical end; \emph{how much} work can be
performed will depend on the forces and vary with these.
Any law which would fix not only the direction of the work,
but also its quantum (``the quantum
satis of the human will'', to speak with Henrik Ibsen), can only indicate an
average or lay down a minimum. Thus it goes
with the ethical laws in their historical form. The standard
that the individual at a given time and in a given people applies
to himself will always bear the stamp of time,
rank and race; he sees himself as he imagines
that an impartial spectator among the human beings he knows
would look upon his conduct. And even if the demands that he
thereby feels made upon himself are satisfied by him more by
the general direction, the animating frame of mind in his
willing, than by a series of adequate actions, it
nevertheless has its \emph{pedagogical} significance that the demands
appear to him bearing the stamp not only of ideality but also
of objectivity. The sooner and the more easily will
the impulse be awakened to get beyond the hitherto existing
state. The contrast between the ideal and the social demand
on the one side, and the individual's actual willing on
the other side, can often be a necessary transitional stage
for the ethical endeavor, as it is possible according to
the individual's own nature, to come about. A spark is necessary to release the bound-up
force. There are natures that first through a painful
% --- p. 61 ---
\setcounter{footnote}{0}%
\opage{61}conflict with the objective norms get their strength sufficiently
freed, for whom there is thus no other way
to harmony than through disharmony. --- The
two propositions: that the ethical law is to formulate the direction
of life, and that it is to work pedagogically, are closely connected.
Both demands are namely satisfied by
the ethical law's setting up \emph{types or models}. A
model indicates the direction and at the same time calls for
imitation. But here a clear
difference does appear between scientific ethics and practical-popular
ethics, between moral philosophy and moral preaching.
Practically it has great significance that the goal be set
high for all, in order that their forces may be released. It
seems indeed to be necessary to demand too much, in order
that what is sufficient may be yielded. Scientific
ethics in its determination of the concept of the ethical law
cannot take regard of practical technique. How easy
or difficult it is in practice to find the law that holds for each
single individual according to its peculiar nature,
makes, scientifically speaking, no difference to the matter. --- On the other
side there is presumably no question either that the more
individualized moral preaching becomes, the more effective
it will be. Moralizing en bloc hardly helps much. It is
thus --- despite the apparent opposition --- nevertheless an
ideal that scientific ethics here brings forth for the
practical art of moralizing.

3. One might say: ``The ethical law is after all a
consequence derived from the principle of welfare, which
cannot be shaken out of regard for the nature of the single individual.
What by virtue of a general principle is right and
proper must be so for \emph{all} individuals.'' But, as
already shown, each \emph{single} individual with its presuppositions
and predispositions belongs nonetheless to the circle of beings
whose welfare is to be promoted. If an inexorable
% --- p. 62 ---
\setcounter{footnote}{0}%
\opage{62}demand is made on it for something that really lies
beyond its capacity, then there is a disproportion, contrary to the very principle of welfare,
between society and individual, a disproportion whose cause does not \emph{need} to lie \emph{solely} with
the individual. The suffering the individual comes to
endure on account of the disproportion between its nature
and the claims of society is therefore not merely to be conceived as
punishment, but also as a fate, just as the suffering
which springs from the lack of agreement between
the individual's nature and the external conditions of nature. When it
turns out that changes in the order of society and in the
conditions of life which are thereby prepared for the single individuals,
often have far greater influence on the number of crimes
than reforms of the penal system\footnote[1]{G. \emph{Tardé}: La criminalité comparée. p. 61: Changez les conditions, s'il se peut, de la société, bien plutôt que son système de pénalité, et sa criminalité se modifiera. --- On the social causes of crime see \emph{Mischler} in ``Handbuch des Gefängsnisswesens'', herausg. von Holtzendorff und Jagemann II. p. 447. 481 ff.}, then
it is clear that one cannot set the historically given
society with its claims as absolutely justified on the
one side, and the single individual with its determinate nature
and dispositions as absolutely obligated on the other
side. The single individual is, after all, in relation to society never
merely cause, but always at the same time effect, and --- what
must especially be emphasized --- it is effect before it is cause.
Even if there lay in the individual's nature unfortunate, antisocial
dispositions, then these might perhaps through a better
order of society be inhibited, diverted or converted
into harmless forms, whereas the present order
perhaps directly affords them conditions for a luxuriant
development. And the more society's order comes about according to
% --- p. 63 ---
\setcounter{footnote}{0}%
\opage{63}conscious deliberation, the greater becomes its responsibility
toward the individual.

If this conception is correct, one must expect to find traces of it
also in theology and jurisprudence, although these proceed from the opposition between an absolute
objective law on the one side and the individual, who has to
obey it, on the other. I find in the role
that grace and pardon play in theology and jurisprudence
a trace of the thought of the individualization of the true law
that I have here sought to develop. In grace and
pardon there is, namely, expressed more or less
clearly a regard for the individuality of the agent and
for his subjective conditions. The orthodox doctrine
even assumes, after all, that the natural man \emph{cannot}
fulfill the law! Grace thus fails to acknowledge the law
by maintaining that the whole individuality of man is not touched
by the opposition to the commands of the law. In the very possibility of
a reconciliation there lies the implication that the law was not absolutely justified
toward the individual; --- otherwise the reconciliation would be a
plain self-contradiction. And when, according to the orthodox
doctrine of reconciliation, God himself, the giver of the law, makes the
sacrifice by which the transgression of the law is atoned for, one might
find in this a symbolic expression of the fact that it is the law, not
man, that was in the real wrong. --- According to the
juridical conception pardon is quite simply ``justice's
self-correction''\footnote[1]{\emph{Jhering}: Der Zweck im Recht. I. p. 428.}, or, as an older jurist\footnote[2]{\emph{Anselm v. Feuerbach}: Aktenmässige Darstellung merkwürdiger Verbrechen. Giessen 1828. I. p. 353.}
puts it more fully: ``To remedy the imperfections of the laws
in their application, wisely to even out the disproportion between the lawful punishment and the culpability of a single
transgressor, between the unchangeable, \opage{64}general law's inflexible severity and the manifold changeabilities of individual
guilt, and thus \emph{to reconcile
justice with equity}, this power
is placed with the right of pardon in the sovereign's hand.''
If justice itself needs to be corrected,
this must come from its imperfection! --- In the
juridical-theological opposition between sin and grace,
crime and pardon, there is an oscillation to both sides
of the center point in which the real truth lies.
Because of the difficulty of hitting the center itself,
one first moves a little beyond it
to one side, then a little to the opposite
side. Instead of moving along the diagonal, one
moves first along the one, then along the other
of the parallelogram's two sides. And as already indicated,
there seem to be natures that can develop only
through such oppositions, and that cannot move
along the diagonal.

In the modern development of the penal system I find
the same thought, more or less consciously, lying at the
basis. --- When the minimum of outward morality that the state
demands of its members in reality makes different
demands of individuals according to their different
nature, then it is the duty of the state to take these
differences into account and to find starting points in the nature of the individuals
for a development toward what the law demands.
Not everything is accomplished by laying down the law and imposing
the punishment. The recognition also seems now to have
awakened that the execution of punishment is of no less importance
than penal legislation and penal sentences, and that
the experiences gathered during the execution of punishment
ought to become of essential significance for the criminal law of the future.
It is precisely in the execution of punishment that the individualizing moment emerges, which penal law and penal sentence can only
% --- p. 65 ---
\setcounter{footnote}{0}%
\opage{65}most imperfectly allow to come into its own.
The experiences gained in the execution of punishment will
therefore react more and more on the ordering of the criminal law.
Thus \emph{Franz von Holtzendorff}
has recently stated: ``The modern theories of punishment must far rather
be derived by way of induction and \emph{from the experiences of the execution of punishment} than, conversely, the execution of punishment from a priori
theories''\footnote[1]{Handbuch des Gefängnisswesens. Hamburg 1888. I. p. 387. --- The following remarks on the development of the prison system are also taken from this work.}. It is also recognized more and more that
what one has to deal with in the penal system is not crimes as objective facts, but the personalities of the criminals.
It has therefore been demanded that the unalterable sentences be
replaced by sentences that fix an indeterminate term of punishment,
so that account may be taken of the individual's character, as
this shows itself during the execution of punishment. A
definite classification of criminals is demanded, with different
treatment according to their different character. In the last century (since the time of John
Howard) the thought has on the whole asserted itself in all reforms of the penal system that the criminal shall, as far as possible, return
to human society not merely cowed and
coerced, but truly improved. The experience has become
generally recognized that, given their
inner and outer circumstances, crime lies infinitely near for many, while
for others it lies very far off. The right of the individual
against the abstract, impersonal law is thus
recognized more and more. Individualization and
humanization go hand in hand and are in reality
inseparable.

In an ideal state only that would be demanded of each which
corresponded to his abilities and powers. The demand would
be adjusted to his peculiarity, just as in
% --- p. 66 ---
\setcounter{footnote}{0}%
\opage{66}education we adjust the demands we make of the child to
its stage of development, so that the demand may give occasion
for the unfolding of the powers and not for the subjugation
of the whole being of the individual. The social disadvantages that would
result from the same demands not being made of all
individuals would then have to be remedied by special measures.
As long as this ideal of a state that educates through all
its institutions and laws cannot be realized,
the state must, through the execution of punishment, perform a part
of the work that it has not been able to carry out in a more direct ethical way.
The actual order leads
to oscillation around the true, ideal order as an invisible
center in hard, painful oppositions.

4. By taking the individual factor into account
the casuistic problems will be solved as if of themselves. Purely
objectively it may be impossible in the particular case to
decide whether one should work rather for a larger society
or for a smaller, whether the main emphasis should be laid
on the dissemination of the culture already won to as many
as possible or on producing new culture, which for the time being
will benefit only a few, etc. But such
questions will as a rule find their decision
in the fact that each person has his abilities and dispositions, which condition
his calling. In the concept of a calling it is expressed precisely and clearly
that a person's duty is determined not merely by the
general ethical principles and the given factual circumstances,
but also by his own individual nature. The concept
of a calling rests on the thought that mastery shows
itself in limitation, since all willing and acting must aim
at something definite, and that the limitation depends essentially
on the individual abilities and powers. It is
these that ``call,'' in that power has a tendency to
awaken impulse. --- If we take activity in the service of culture
as an example, there will be those whose abilities
% --- p. 67 ---
\setcounter{footnote}{0}%
\opage{67}best make them fit to work in a narrow circle,
and whose activity therefore naturally becomes limited
to providing for their own maintenance and that of those nearest to them,
so that their work at most can have significance only indirectly
beyond the narrower circle. Others
there will be who are suited to work upon and for many,
and who will therefore be good instruments for spreading
the culture won so far still further. Still others will have
instinct and power to strike out on new paths, to
try new possibilities. Despite contradiction and lack of
understanding on the part of the multitude they carry on their work
and tend the first sprouts of the culture of the future. For the individual who knows himself there will
here be no doubt as to where he should turn. But the main difficulty
is precisely that of knowing oneself, as
already \emph{Socrates}, the founder of ethics, impressed upon us.
The Socratic fundamental thought shows itself precisely through the conception
I here seek to maintain, of the law as
individual, in its whole great and deep significance. Self-deception
is so easily possible where it is a matter of finding out
one's own ability and marking out one's goal and thereby one's limit.
No objective criterion is wholly adequate here, and something
must always be ventured, one must act in reliance on what the
personal conviction proclaims. Every choice is
here a practical hypothesis, the verification of which it is the task
of life to supply. If the individual in his ethical formula
inserts incorrect coefficients, the task is not solved.

The same holds when, instead of the work of one's calling,
we take \emph{self-control} as an example to
illuminate the law of relativity in ethics.

Self-control in particular spheres can occur
together with a complete lack of self-control in
other spheres. The Indians, e.g., are able to
endure the most terrible tortures when they fall into enemy
% --- p. 68 ---
\setcounter{footnote}{0}%
\opage{68}hands, or when by imposing on themselves the most
exquisite sufferings they wish to atone to and honor their gods. Illness, too,
they bear without complaint, and they demand of their
women that they give birth without cries of pain. In their
outward bearing they always preserve the greatest calm, decorum,
and seriousness, even when the most violent passions rage in their
inner being. But in the face of enjoyments such as food, drink, and games
they cannot govern themselves. Only with great exertion
did some chiefs succeed in getting it accepted that drunkenness
should be regarded as a vice. --- This shows that the
capacity for self-control is not a constant factor,
independent of whatever else stirs in the person.

Still more clearly does this appear in self-control
in the sphere of the sexual instinct. The different individuals
are here most differently equipped. There are those in
whom the instinct works in harmony with sympathetic feelings,
so that it demands no satisfaction except where it
forms the natural basis for heartfelt devotion to
another individual, who is embraced with joy and admiration.
This is of all the greater importance because the satisfaction
of the sexual instinct falls under ethical consideration chiefly for the reason
that it is not possible except by means of another individual. Apart from this
relation and the responsibility arising from it the sexual instinct would
be to be judged only from a purely individualistic
or egoistic standpoint, that is, according to hygienic
considerations or according to regard for the harmonious development
of the individual's mental and bodily faculties as a whole.
The great problem, the ethical difficulty, arises
precisely from the question how this mighty instinct, which sets
such strong feelings in motion and is so significant
for the welfare of the single individual himself, is to be
satisfied in such a way that it does not cause physical and moral
suffering to other individuals and does not intrude disturbingly into
% --- p. 69 ---
\setcounter{footnote}{0}%
\opage{69}the right social order. --- Besides the previously mentioned
happily constituted individuals there are others, in whom the
instinct stirs in all its elemental brutality, as
a purely physical urge, in the satisfaction of which the object
itself may be a matter of indifference. --- In still others the instinct is
comparatively strong, but they are nevertheless able to divert attention
from it by ideal pursuits or other work.
--- In still others it cannot be mastered in this
way. There are on the whole countless gradations here, right
up to the ``morbidly heightened sexual impulse,'' which according to Krafft-\emph{Ebing}
is frequently found in individuals with a ``neuropathic
constitution,'' individuals who ``for a great part of their lives
suffer under the burden of this constitutional anomaly in
their life of the impulses''\footnote[1]{Psychopathia sexualis. 2. Aufl. p. 34.}. --- The two sexes are as a rule differently
disposed in this respect. The harmony between the sexual instinct
and ideal devotion to another individual
is certainly far more frequent in woman than in man,
which may well be connected with the fact that the male throughout the animal world
is the more active and aggressive party in reproduction,
which is why the instinct in him has such an elemental
and ruthless strength as it only rarely has in
the female. --- But in men, too, the variation of the strength of the instinct
and of its relation to other motives is, as said, infinite. In the dispute that is frequently conducted about sexual
morality, these countless individual
differences are as a rule overlooked, which are nevertheless certainly the most important
element in the whole matter. From the opposite
sides equally absolute assertions are put forward. From the one
side strict, universally valid commands are formulated without
seeing that one thereby lays a most different burden
on the different shoulders. From the other side one postulates
a physiological necessity to which everyone must
% --- p. 70 ---
\setcounter{footnote}{0}%
\opage{70}succumb to. This is a fate that does not exist.
The sexual instinct appears under other conditions in man
than in the animal. When it awakens in man,
the capacity for self-control has already developed in other spheres.
Ideas and feelings have been set in motion
that do not permit the direct dominion of the instinct,
and that, when attention and interest in
ideal self-maintenance are aroused, can make it possible
to subdue it and bring its manifestations into harmony with the other
demands of life. Only it must not be forgotten that self-control
as a purely negative virtue is a psychological impossibility.
It is a proposition that is all too often forgotten in ethics,
that an impulse can be inhibited only by another impulse. In self-control it therefore
depends not merely on how strong
the inclination is that is to be subdued, but also on whether
there is the possibility present of other inclinations that
can lay claim to the mental energy. What matters
is that there be a great love, a great enthusiasm
that can fill the soul. The law of relativity in the sphere of ethics
proves here to be a simple consequence of the law
of the conservation of energy\footnote[1]{One could here also make use of the law of inertia. \emph{Spinoza} proves his proposition: affectus coërceri nec tolli potest, nisi per affectum contrarium et fortiorem (Eth. IV, 7) from the law of inertia.}. The energy that our consciousness
or our brain commands is at every time limited;
the more of it that is consumed for the satisfaction of one impulse,
the less remains for the other impulses. Annihilation
is impossible, but conversion into other forms
may be possible. The great heroes of self-control possessed
a great enthusiasm for other sides of life than
those they turned away from. The great self-control that the Indians
show in the face of wounds and pain in war, captivity,
or religious worship is conditioned by their concepts of honor,
% --- p. 71 ---
\setcounter{footnote}{0}%
\opage{71}their tribal feeling, their religious faith. In the face of sensual
excesses their circumstances of life had given no occasion
for the arising of counterbalancing motives. As \emph{Leslie
Stephen} has remarked, there develops in man,
as he advances in intellectual, aesthetic,
and sympathetic respects, a higher sensibility, which brings with it
a feeling of disgust and contempt for sensual
excesses, apart from their calculable consequences, a sensibility
that can act just as directly and imperatively
as any of the elementary sense faculties\footnote[1]{Science of Ethics. p. 102. 195. 200. --- See above p. 23.}.

The struggle of self-control goes on until the new
employment of energy has gained complete preponderance
over the earlier. This struggle becomes of a most different
kind in different individuals, according to the relation between
the original or acquired inclinations. In
some the conversion takes place even without struggle; but there
are those who struggle hard without gaining firm ground under
their feet. The work of will that goes on in these latter
is far more serious than that which is performed by the former, although the
outward result is far poorer. When Jesus preferred
to associate with harlots rather than with Pharisees, it was perhaps
because he found in them an inner striving and longing which,
though unable to overcome the opposing barriers, yet
stood in closer kinship with the ideal ethical law
than the self-satisfied correctness of the Pharisees. ---
--- The law of relativity throws light on many ethical questions
that would otherwise be insoluble. It teaches us to
unite the ideality of ethics with the realistic recognition
of the multiplicity of life, in that it shows us that the real
ethical development proceeds from a multitude of different
starting points, from very different levels
toward one and the same goal.

% --- p. 72 ---
\setcounter{footnote}{0}%
\opage{72}5. Where the task arising from the general principles and the
given situation fully agrees
with the individual's own ability and impulse, we have the happiest
--- and when the task is great --- the greatest people.
What is characteristic of these is the inner security
and harmony with which they --- despite all inner and outer
struggles --- advance toward their goal and fill their
place. Here the ``organic morality'' so called by \emph{Herbert Spencer}
is present. That there are such
natures is a fact that cannot be denied, however
one may then wish to explain it. It is they who provide models
for others. It must here be maintained, however, that
the task is set the individual not merely out of unchangeable
situations given once for all, but that the individual, by
virtue of his ability and impulse and in accordance with the highest
principles, can also precisely set himself the task of
changing the situation at hand, for which this of course
must itself contain a possibility. In the
great pioneers there develop abilities and impulses that
cannot find their satisfaction under the given
conditions, and that therefore themselves produce new
conditions out of these. Such individuals probably arise through
one of those variations that according to Darwin's hypothesis are
the precondition for a new form of life to be able to form.
Their right coincides with their ability. By the very
unfolding of what lies in their nature they provide
models for a new direction in one of the domains of human life
and thereby become means for the progress of the race.
Nature and art, nature and duty act here immediately
together. They bear the law within themselves in a higher sense
than can be said of those in whom no new possibilities
for a higher development lie. Yet ``organic morality''
may also be found in those whose tasks are essentially determined by the conditions
as they are given,
% --- p. 73 ---
\setcounter{footnote}{0}%
\opage{73}if their nature leads them immediately both to the
activity and to the limitation and self-control
that the conditions require.

The complete opposite of such natures
would be present if there were to be individuals
in whom absolutely no movement
or striving in the direction of the ethical ideal, not even the faintest, could be released.
All the coefficients in the ethical formula would then be
zero. No demand or appeal, however pedagogically formulated, would
be able to act as a stimulus. Here we come, among other things, to the question whether there are incorrigible
criminals, a question that in the most recent times
has been answered in the affirmative with great decisiveness by
\emph{the Italian school of criminal psychology and philosophy of law}.

This school has great merits in respect of science.
The usual ethical and juridical mode of consideration
stops at the action and the motives lying nearest
and does not go back to the causes that have
produced the whole character of the agent. The formerly
prevailing indeterminism has favored this conception.
Criminal anthropology and psychology, which
owe so much to the Italian school, are working to
remedy this one-sidedness. It seeks to establish
the criminal's bodily and mental peculiarities
and to demonstrate the causes of these. On the basis of these investigations it has
set up the criminal type
(homo delinquens) as a special variation of the
general human type. Among the bodily peculiarities
of this type \emph{Lombroso} mentions in particular such
properties of the cranium, brain, and physiognomy
as must be regarded as deviations from the
type of \emph{civilized} man (\emph{regressive} deviations), among others small
cranial capacity, a small and narrow forehead. Among the
% --- p. 74 ---
\setcounter{footnote}{0}%
\opage{74}mental peculiarities he emphasizes a remarkable insusceptibility
to pain, with which presumably a no less
remarkable lack of compassion stands in natural connection.
The vasomotor reflexes are weak on the whole,
and it is rare for the true criminal natures to be able
to blush. The moral feeling is lacking, which (as
\emph{Despine} had already stated) shows itself in the great
rarity of repentance. This whole type is now, according to Lombroso
and the Italian school, to be explained by atavism.
It arises nowadays by deep strata from the
prehistory of the human race pushing up again in the midst of civilization.
In no other way, it is thought, can one
explain the unmotivated and ruthless character of the conduct of the true
criminals, just as, too, the frequency of relapses
and the slight effect of the best-devised systems of punishment
are understood, when it is forces from primeval times that
stir anew.

The moral sense, says a prominent Italian
scholar, Garofalo (in his ``Crimonologie''), is the work of centuries. If now --- as experience shows ---
some individuals have not become partakers in this work,
but suffer from an innate moral insensibility, which
shows itself in compassion's not stirring in them,
not even in its most elementary forms --- how
can one then imagine that any influence
and education whatever in the course of the single individual's lifetime
should be able to accomplish what for the rest of the race has
cost ever-repeated experiences over the course of millennia?
Man is good by instinct, not by reflection; when
then the instinct is lacking, what other force could then
remedy it? Individuals who do not (like
occasional criminals) offend on account of passing
occasions, but who again and again return
to their criminal ways, or who through the comparatively slight motives of the
% --- p. 75 ---
\setcounter{footnote}{0}%
\opage{75}action and the revolting manner
in which it is carried out give insight into an ingrained antisocial
disposition, --- such individuals it is impossible
for a civilized society to assimilate. They would
fit into far earlier civilizations, and nowadays
perhaps in Dahomey or on the Fiji Islands; but
it is hopeless to want to make them elements in our
present social life. Toward them society can then
have no duty. On the other hand it has toward itself the
duty of eliminating such ``molecules'' as would harm
life in the remaining organism. It is an individualistic
one-sidedness to want to maintain that criminal individuals of
the kind described should have rights as against society.
Social necessity is unconditional, and the
social reaction tolerates no limitation.

So far \emph{Lombroso} and Garofalo. --- The view maintained by these
scholars interests us here insofar
as it assumes an absolute moral insensibility
in some individuals, who would thereby be excluded from
belonging in any degree whatever to the multiplicity of
starting points from which human development
moves toward its ideal. Not only individual
relativity, but also psychological relativity
would fall away. We would have before us, in the disposition and mode of thought of the true criminal natures, a system
that was fundamentally opposed to the ethical system, on which
human development, consciously or unconsciously, is
built. And there should be absolutely no possibility for
that system to be overcome by this one except by a
war of extermination. Before we accept this assumption,
we must subject it to a closer discussion.

a. It is always uncertain to infer from particular actions
to a person's innermost character and the possibilities
it contains. A person's character and its
% --- p. 76 ---
\setcounter{footnote}{0}%
\opage{76}possibilities are not exhausted by the single action. According to Garofalo the crime need not even
be especially great, and yet the examination of the criminal's
character, by means of the information that can be
procured, is supposed to be sufficient for a judge trained in anthropology
and psychology to decide
whether the individual is among those who can accommodate themselves to
the demands of civilized society, or whether he must be ``eliminated.''
Indeed, not even an action is needed; an attempt is supposed
to be sufficient!

Just as uncertain is the inference that is drawn from
the absence of repentance. It has been established especially in
condemned criminals or in prisoners such as those with whom Dostojevski
had to share prison in Siberia, and who just
were not the object of any truly pedagogical treatment.
Because repentance cannot be awakened in the short
time between the commission of the crime and the execution of the death sentence,
and because it could not be awakened in criminals
who were subjected to the mode of treatment that until a short
time ago was the usual one in prisons,
one cannot nevertheless infer from this that it \emph{under no
circumstances} could be awakened. --- The frequency of relapses
likewise points to ingrained, but not necessarily to
incorrigible, inclinations.

In the case of imprisonment Garofalo will have no
definite time fixed by the sentence, among other reasons ``because
one cannot measure the share in the crime that the influence of
the surrounding conditions had, and that which the individual's
anomaly itself had.'' Only after a longer period of observation,
in which one can come to know the individual's character,
can the limit of the term of punishment be
fixed\footnote[1]{La Criminologie. p. 395.}. But how can he then in some cases
% --- p. 77 ---
\setcounter{footnote}{0}%
\opage{77}demand absolute elimination? The uncertainty in the knowledge
of the criminal's character, which means that the length of the imprisonment
cannot be fixed in advance, must
also make it impossible to establish absolute elimination
as the only way out!

b. When the incorrigible antisocial dispositions are explained
by atavism, a new problem arises, namely
how such marked peculiarities, which concern not
merely a single side of the individual's character, but
the whole character, have been able to maintain themselves within the race
alongside peculiarities of an entirely opposite kind.
How has the thread through which the single
individual is connected with the primeval generations been able to avoid
being spun together with the thread by which
he is connected with the nearest generations and with
the institutions and traditions that these have produced
and handed down? In other words: even if the explanation of
the criminal type by atavism is correct\footnote[1]{G. \emph{Tarde} has in his ``Criminalité comparée'' criticized this explanation and arrives at the result that, even if one must acknowledge unquestionable physiological and anatomical analogies between the born criminals and the savages, and even if the former --- like all monstra --- present regressive traits, yet these are found in them \emph{combined in a peculiar way}, and one must not judge our forefathers by this sample. (p. 46). But then it is a matter of explaining ``the combination'' --- and that cannot, after all, be independent of social development.}, so \emph{atavism
is no proof} that \emph{society has no guilt
and therefore no duty toward the individual either}.
For when the antisocial germs have been able to maintain
themselves so long in latent form within the race and have not
been transformed, this must nevertheless come from deficiencies
and imperfections in the life of the race. The perfection of social conditions
and of the ethical life generally must
nevertheless be measured, not only by what positive progress
% --- p. 78 ---
\setcounter{footnote}{0}%
\opage{78}is made, but also by what unfortunate possibilities
are crowded out and abolished. The more perfect a social order
is, the more it acts on the innermost being of the single individuals.
Atavism is thus itself a sign of
social imperfection and does not justify society's
being placed over against the criminal individual as the absolute
right against absolute lawlessness. --- Strangely enough
the Italian school, in that it places the criminal
individuals outside society both in the physical and
the ethical respect, arrives at the same result as indeterminism
and individualism, though these are otherwise its antipodes.

c. When Garofalo, like \emph{Despine} before him
(and \emph{Anselm von Feuerbach} still earlier), teaches
that moral insensibility can remain latent until
external occasions betray it, --- when he teaches that even
if poverty is not the real cause of crime,
yet many people would have remained
latent criminals if they had not been poor,
and if they had lived in a different milieu, --- then
it seems clear that society must be a contributory
cause, since its order always more or less conditions
the milieu in which the single individuals develop.
One may distinguish ever so sharply between the mere occasion
and the real cause; this distinction will always
be more or less artificial, in that \emph{the complete
cause consists of both the external occasions
and the inner dispositions}. There is so much the more reason
to emphasize this, since it is after all the ``occasions'' that
bring about the enormous difference between two latent criminals,
that the one is declared unassimilable, while the
other undisturbed continues his life as a molecule of society!


d. \emph{Inherited dispositions always have a certain
indefiniteness}, which makes it possible that by education
% --- p. 79 ---
\setcounter{footnote}{0}%
\opage{79}and external conditions they can be developed in different directions.
What is inherited is only dispositions. It cannot
be proved that there is present a fatalistic determination by nature
that would make all practical experience, lived experiences, and exertions
meaningless. Even if a good education and
fortunate conditions should not be able to stifle the evil
germ, but at most could bring it about that it unfolds
without social disadvantages, a bad education will nevertheless
do an extraordinary amount toward developing the evil germ.
One will therefore not be able to draw any absolute line
between what springs from the individual in and of himself, and
what is due to the conditions in which he is placed, and which
for the most part depend on the order of society. --- Very
often there is no absolute egoism present. In criminal natures
the family feelings can often be very
strong, so that the criminal instincts manifest themselves only
toward strangers. Here there might perhaps lie a
possibility for the awakening of more far-reaching sympathy, just
as it is from the family feelings that love of mankind
has from the first developed to embrace larger circles.
Another starting point might be sought in the criminal's self-esteem
and urge to assert his person. Dostojevski
remarks that the worst punishment for the criminal
would be \emph{useless} work. This is certainly connected
with the fact that useless work, like pouring into the vessel of the Danaids, leaves the
personality standing as a nothing, which has significance neither for itself
nor for others. This last remnant of self-esteem would, without
influence from outside, not be enough to overcome the laziness that
is a frequent part of the criminal nature; but it can furnish a starting point for education, and the modern systems of punishment are
surely with justice convinced that it is above all through
work that education is to take place.

e. Because our institutions have not yet solved
% --- p. 80 ---
\setcounter{footnote}{0}%
\opage{80}the great task of preventing the increase of crimes, especially of relapses,
one cannot therefore infer that there really should be
absolutely incorrigible criminals,
individuals in whom no ethical striving whatever should
be able to be awakened, and who could not at all be brought
into relation to the ethical ideal. Even if a complete
return to human social life should not
be possible, an impulse could still be awakened that
led forward in the direction of the ideal, were it only with
small steps. A reconciliation would then be possible, and the
\emph{absolute} elimination not necessary. --- There is no
reason to believe that the development of the penal system should
have reached its peak. The reform is only at its very first
beginning. People have hitherto been more inventive in
devising punishments that could satisfy the urge for retribution
and the desire to deter than in finding
means and ways of influencing the character of criminals
in order, through this belated education, to restore
what through the order of society had been neglected at
an earlier stage. This came from the external way in which
crime was formerly regarded. People held \emph{to the fact that} it had been committed, and laid less weight on
\emph{how} and \emph{why} it was committed. The Italian
school, through its works, which illuminate the character of the criminal
from so many sides, will contribute much to putting a complete
end to this incorrect mode of consideration.
This merit is independent of the validity of the
conclusions that the school itself draws from its investigations.

After all this it therefore still stands as an open
question whether there are human individuals who
are altogether incorrigible, so that the ethical law, however
individualized it might be made, should not be able to
come to assert itself in them. For the general
theory that I have wished to maintain, the answer to this
% --- p. 81 ---
\setcounter{footnote}{0}%
\opage{81}question is not of decisive importance. But
if my theory is correct, it calls for great caution
before one declares individuals to be \emph{absolutely} insusceptible
in the ethical respect. If each individual has
his own individual law, that question is far more difficult
to answer than if one had a right to proceed
from an objective law that should be capable of being fulfilled by all.
What is demanded is a sense for individual differences and
a patience that gives itself time until the hard
shell can be broken apart. Toward criminals (``the criminal type'')
people have long acted as Europeans
are wont to act toward the other races.
Just as these have summarily been declared to
lack the capacity for civilization\footnote[1]{The considerations with which Th. \emph{Waitz} counters the assertion of the Indians' lack of capacity for civilization are of interest in this connection. He says, among other things: ``The Indians' long standing still at a very low stage of culture proves little or nothing against their endowment, as long as they have not yet for any longer time been placed in such conditions as at all make progress possible, and this has so far not happened. Indeed, still more. When for many centuries, as far as we know, nothing has acted upon a people from which one could promise oneself a transforming influence on their organization and their mode of life, on their outer and inner life, then it is probable that far stronger powers and a continuous influence sustained over a far longer time are needed in order to bring about any transformation of them at all, however slight, than in the opposite case.'' Die Indianer Nordamerikas. Leipzig 1865. p. 74.} and thereby doomed to
destruction, so those have been declared to lack
moral capacity and thereby doomed to ``elimination.''

6. In conclusion I will develop somewhat further the thought
expressed in the foregoing, that the significance of the individual differences
in ethics is not exhausted by their
entailing the necessity of a thoroughgoing
% --- p. 82 ---
\setcounter{footnote}{0}%
\opage{82}individualization of the ethical law, or by life's thereby gaining
greater fullness and richness than if it were to and could be
cut to fit certain definite, general rules. The recognition
of their significance also makes the advance of the race
intelligible. --- Development within the organic
world becomes possible, as Darwin has taught, only
through an uninterrupted varying taking place. Of the
``spontaneous'' individual variations, those are selected and retained
that are most serviceable for the continued existence of life.
These variations thus retained become, accordingly, typical
for the continued development of life. We have an analogy to this
in the ethical domain. Those in whom the traits of character that the ethical law demands
develop with the greatest power and to the highest degree, and by whom the
inner and outer resistance is overcome with ease or,
even if it happens through hard struggle, in such a
way that they thereby break through barriers that have hitherto
halted development, --- they do not, to be sure, in a certain sense
stand ethically higher than the many who must struggle hard
merely to attain an average level of development, and who essentially
keep to the hitherto beaten paths, but they nevertheless stand
as models who guide, encourage, and console; they indicate
not merely \emph{the way and the direction}, as patterns for imitation,
but show also \emph{what possibilities are contained in human} nature. Despite the
great individual differences all human beings are nevertheless akin, and it is a natural inference
that what those furthest advanced have attained must also to a certain
degree be capable of becoming accessible to the others. As
\emph{Herbert Spencer} has said: ``What now characterizes
the exceptionally great may be expected some day to
characterize all. For what the best human
nature is capable of lies within the reach of the whole
% --- p. 83 ---
\setcounter{footnote}{0}%
\opage{83}human nature''\footnote[1]{Data of Ethics. p. 256 f.}. History shows that the people
who stand as ethical models for long ages were
not always such as had an inner being full of strong and
lasting inner oppositions, or whose character bore traces
of great inner struggles, but rather harmonious natures,
who lived and breathed in great thoughts, to which the
other parts of their nature bowed almost without resistance.
And this is natural enough, for he whose power
is claimed by the inner struggle will not easily be able to present to others
the picture of a vigorous personal activity directed toward a significant
goal. In the religious idea this has even been carried through to such a degree
that the model has been made sinless and people have wished to remove every
trace of real, human struggle and striving.
This is to shoot past the mark; but it bears witness
to the urge for exalted models, an urge that
cannot fall away if life is not to become flat and dull.

The solidarity of the race is thus not canceled by the
individual relativity of the ethical law. The unity of the race consists
precisely in and through the multiplicity of individuals.
It is, as I have sought to show, in the common direction
and tendency of life that the justification for speaking of ethical law
is grounded. And thereby the unity of the race is also established
in spite of all differences.
\end{document}
